% anthology.tex — content-style disciplined (~1 box/ch, prose spine)
% Generated by gen_anthology.py — follows references/content-style.md
% ============================================================================
% preamble.tex — tested LuaLaTeX preamble for Hebrew + math + figures PDFs (v3.23)
% Battle-tested: compiled an 80-page Hebrew quantum-computing summary AND a
% 197-page Hebrew computer-architecture anthology (TikZ diagrams, code listings,
% colored boxes) — both with 0 missing characters.
%
% Usage in your document:
%   % ============================================================================
% preamble.tex — tested LuaLaTeX preamble for Hebrew + math + figures PDFs (v3.23)
% Battle-tested: compiled an 80-page Hebrew quantum-computing summary AND a
% 197-page Hebrew computer-architecture anthology (TikZ diagrams, code listings,
% colored boxes) — both with 0 missing characters.
%
% Usage in your document:
%   % ============================================================================
% preamble.tex — tested LuaLaTeX preamble for Hebrew + math + figures PDFs (v3.23)
% Battle-tested: compiled an 80-page Hebrew quantum-computing summary AND a
% 197-page Hebrew computer-architecture anthology (TikZ diagrams, code listings,
% colored boxes) — both with 0 missing characters.
%
% Usage in your document:
%   \input{preamble.tex}     % (or paste these lines at the top)
%   \begin{document}  ...  \end{document}
%
% Compile with:  lualatex -interaction=nonstopmode doc.tex   (twice, for TOC)
% Requires the Hebrew fonts installed (run scripts/setup_fonts.sh first).
%
% SHORT doc (exam + solution): change `report` -> `article`.
% STUDY doc meant to be read/annotated: 12pt + \linespread{1.3} below are
%   deliberate readability choices (NOT padding) — keep or trim to taste.
% ============================================================================
\documentclass[12pt]{report}          % 12pt: readable for a study guide; 11pt is fine too
\usepackage[no-math]{fontspec}        % (1) ★ no-math — THE critical line (see recipe.md)
\usepackage[bidi=basic]{babel}        % (2) RTL handling built into the engine
\babelprovide[import,main]{hebrew}    % (3) Hebrew = main language, default RTL
\babelprovide[import]{english}        % (4) English secondary, for correct LTR runs
% Latin fallback registered BEFORE \babelfont so a lone English word in Hebrew
% prose does not tofu. Fonts resolved by FAMILY NAME (portable; no hard path).
\directlua{luaotfload.add_fallback("hebfb",{"DejaVuSerif:mode=node;"})}
% v3.22: Frank Ruhl Libre ships NO italic face. Instead of silently-upright
% \emph (v<=3.20) or bold-as-emphasis (v3.21), synthesize a slanted face with
% FakeSlant — real, VISIBLE italic emphasis with correct \emph nesting
% semantics (inner \emph toggles back upright). Verified empirically at
% 300dpi: slant clearly visible, 0 errors, 0 missing glyphs.
\babelfont{rm}[RawFeature={fallback=hebfb},
  ItalicFont={Frank Ruhl Libre}, ItalicFeatures={FakeSlant=0.22},
  BoldItalicFont={Frank Ruhl Libre Bold}, BoldItalicFeatures={FakeSlant=0.22}
  ]{Frank Ruhl Libre}
\babelfont{sf}[RawFeature={fallback=hebfb}]{Heebo}
\babelfont{tt}[RawFeature={fallback=hebfb}]{DejaVu Sans Mono}

\usepackage{amsmath,amssymb,mathtools}
\usepackage[margin=2.3cm,headheight=14pt]{geometry}
\linespread{1.3}                      % readable line spacing for study/annotation docs
\usepackage{xcolor}
\usepackage{mdframed}
\usepackage{listings}                 % code listings (assembly / C / pseudo-code)
\usepackage{tikz}
\usetikzlibrary{arrows.meta,positioning,calc,shapes.geometric,fit,backgrounds}
\usepackage{pgfplots}\pgfplotsset{compat=1.18}
\usepackage{quantikz}                 % quantum circuits (remove if unused)
\usepackage{booktabs}
\usepackage{longtable}
\usepackage{array}
\usepackage{float}                    % [H] = place table exactly here (reference doc: no drifting)
\usepackage{enumitem}
\usepackage{needspace}

% ============================================================================
% ★★ v3.4 — THE DIRECTION-ISOLATION TRIO (the heart of mixing Hebrew + Latin)
% Mental model:
%   \h{...}    Hebrew (RTL) island      — Hebrew text inside an LTR context
%                                          (TikZ nodes, the english `tikzpic`).
%   \en{...}   English (LTR) isolate    — Latin text / parentheses inside a
%                                          Hebrew (RTL) context. KEEPS PARENS
%                                          AND ORDER CORRECT.
%   \code{...} \en + monospace          — code inside Hebrew / inside a frametitle.
% CRITICAL: never wrap Latin in \h{} (it is for HEBREW). Inside a TikZ node,
% Latin in \h{} renders REVERSED (Datapath->htapataD, CPU->UPC) and parens
% mirror. Isolate every Latin/paren chunk with \en{} or \code{}. See
% references/figures-boxes-listings.md.
% ---- TEXT INSIDE MATH (v3.8) ----------------------------------------------
% Put ANY text that sits inside math in a \text{}. The ONLY thing that is
% invalid in math is a bare \h/\en (= \foreignlanguage) used DIRECTLY, outside
% \text -> "\rmfamily invalid in math mode". Inside \text:
%   Hebrew label           -> \text{...}            (bare Hebrew renders fine;
%                              \h is OPTIONAL/harmless: \text{\h{...}} also works
%                              and both render identically + scale in scripts)
%   plain Latin word        -> \text{...}            (no marker needed)
%   Latin WITH ( ) or [ ]   -> \text{\en{...}}       (\en forces LTR; a plain
%                              \text mirrors the parens in RTL -> "Spec (TNR(" )
%   f(Hebrew/text arg)      -> \text{f}(\text{...})  -- keep the parens at MATH
%                              level, NOT \text{f(arg)} (parens inside the box mirror)
%   code in math            -> \mathtt{...}  (or \texttt{}, math-guarded below)
% NEVER write \en{...} or \h{...} straight into $...$ / \[..\] (outside \text);
% always via \text{}. (\mbox{...} also works but does not scale in scripts.)
% ============================================================================
\newcommand{\h}[1]{\foreignlanguage{hebrew}{#1}}
\newcommand{\en}[1]{\foreignlanguage{english}{#1}}

% ============================================================================
% ★★ — fixes for defects that are INVISIBLE in the log and in the checks
% and only surface on a VISUAL RENDER of the PDF. (Each verified empirically.)
% ----------------------------------------------------------------------------
% v3.22: \emph is now REAL slanted emphasis (FakeSlant italic face defined on
% the \babelfont{rm} line above), replacing the v3.21 \emph->\textbf hack: it
% preserves \emph's toggle semantics and keeps bold free for its own role.
% Prefer bold-as-emphasis for your doc? one line: \renewcommand{\emph}[1]{\textbf{#1}}
% page number: force the digits LTR so they never reverse (18 -> 81) in the RTL
%   footer. That reversal is a NON-DETERMINISTIC bidi glitch (observed once in
%   ~80 pp); this LuaLaTeX build is not byte-reproducible, so the defensive
%   LTR wrap is the only reliable guard. Latin digits are always LTR anyway.
\renewcommand{\thepage}{\foreignlanguage{english}{\arabic{page}}}
% emergencystretch: Hebrew does not hyphenate, so an unbreakable LTR island
%   (\en{...} / an inline math run) landing at the margin can leave a small
%   (~2-13pt) overfull with no legal breakpoint. This rescues ONLY those
%   inline-prose tie-breaks; it is a ceiling used only on offending lines and
%   does NOT touch display-math / listing / table overflow (those still surface
%   in the log for the content-level overflow pass). Verified: baseline overfull
%   -> 0 with this set, with no change to pagination or underfull.
\setlength{\emergencystretch}{3em}
% em-dash: --- does NOT become — under Frank Ruhl (the font lacks tlig; verified
%   the width is identical with/without Ligatures=TeX, while the mechanism works
%   on Latin Modern). TYPE THE CHARACTER — (U+2014) DIRECTLY; likewise – (U+2013).
%   check_bidi_figures.py flags a --/--- that is adjacent to Hebrew.
% ============================================================================

% ---- STRUCTURAL DIRECTION-SAFETY for \texttt (v3.7) -------------------------
% A parenthesis, bracket, or $ INSIDE \texttt{} mirrors/migrates in RTL prose
% (MIPS 0($s0)->"(s0$)0", $rs->"rs$", arr[i]->"[arr[i") because \texttt sets the
% font but NOT the direction. Redefining \texttt to LTR-isolate its content fixes
% the WHOLE class at once (parens, brackets, $) so plain \texttt{lw $t0,0($s0)}
% is safe — no need to remember \code{}. Guarded by \ifmmode because
% \foreignlanguage is invalid in math (there, plain monospace is already LTR).
% NB: \texttt{Hebrew} stays tofu — the monospace font has no Hebrew anyway; put
% Hebrew in \h{}/normal text, never in \texttt{}. To disable, delete this block.
\let\origtexttt\texttt
\renewcommand{\texttt}[1]{\ifmmode\origtexttt{#1}\else\foreignlanguage{english}{\origtexttt{#1}}\fi}
\newcommand{\code}[1]{\ifmmode\origtexttt{#1}\else\foreignlanguage{english}{\origtexttt{#1}}\fi}

% ---- a thin decorative rule helper ----
\newcommand{\sep}{\par\vspace{2pt}\noindent\textcolor{cRule}{\rule{\linewidth}{0.6pt}}\par\vspace{4pt}}

% ---- colors ----
\definecolor{cBlue}{HTML}{1F4E79}   \definecolor{cBlueBg}{HTML}{EAF1F8}
\definecolor{cGreen}{HTML}{2E6B4F}  \definecolor{cGreenBg}{HTML}{EAF4EE}
\definecolor{cAmber}{HTML}{9C6A00}  \definecolor{cAmberBg}{HTML}{FBF3E2}
\definecolor{cPurple}{HTML}{5B2C6F} \definecolor{cPurpleBg}{HTML}{F3ECF7}
\definecolor{cTeal}{HTML}{0E6E6E}   \definecolor{cTealBg}{HTML}{E6F4F4}
\definecolor{cRed}{HTML}{8E2B2B}    \definecolor{cRedBg}{HTML}{FBEAEA}
\definecolor{cGray}{HTML}{555555}   \definecolor{cRule}{HTML}{1F4E79}

% ---- math operators ----
\DeclareMathOperator{\Tr}{Tr}
\DeclareMathOperator{\rank}{rank}
\DeclareMathOperator{\diag}{diag}
\DeclareMathOperator{\sgn}{sgn}

% ---- bra-ket (quantikz defines \ket/\bra/\braket; add the rest safely) ----
\providecommand{\ket}[1]{\left|#1\right\rangle}
\providecommand{\bra}[1]{\left\langle#1\right|}
\providecommand{\braket}[2]{\left\langle#1\middle|#2\right\rangle}
\newcommand{\dyad}[1]{\left|#1\right\rangle\!\left\langle#1\right|}
% v3.21: general outer product |a><b| and expectation value <A> — agents reach
% for these by reflex (esp. \expval); quantikz defines neither, so their absence
% threw "Undefined control sequence" in a real build. \providecommand = safe.
\providecommand{\ketbra}[2]{\left|#1\right\rangle\!\left\langle#2\right|}
\providecommand{\expval}[1]{\left\langle#1\right\rangle}
\newcommand{\abs}[1]{\left|#1\right|}
\newcommand{\norm}[1]{\left\lVert#1\right\rVert}
\newcommand{\half}{\tfrac12}
\newcommand{\R}{\mathbb{R}}\newcommand{\C}{\mathbb{C}}\newcommand{\Z}{\mathbb{Z}}
% identity operator. NB: \mathbb{1} silently renders a WRONG glyph (msbm has no
% blackboard digits). Use \mathbb{I} = 𝕀 (works with amssymb, no extra package).
% Do NOT "fix" this back to \mathbb{1}.
\newcommand{\id}{\mathbb{I}}

% ============================================================================
% Colored boxes — mdframed (NOT tcolorbox, which breaks under bidi).
% ★★ v3.4: innerleftmargin/innerrightmargin give code listings room so their
% frame/line-numbers do not spill out the box border (see listing styles below).
% ============================================================================
\newmdenv[linecolor=cBlue,backgroundcolor=cBlueBg,linewidth=1.1pt,
  frametitlebackgroundcolor=cBlue,frametitlefontcolor=white,
  innertopmargin=6pt,innerbottommargin=7pt,skipabove=9pt,skipbelow=9pt,
  innerleftmargin=9pt,innerrightmargin=9pt,roundcorner=3pt]{defbox}   % הגדרה
\newmdenv[linecolor=cGreen,backgroundcolor=cGreenBg,linewidth=1.1pt,
  frametitlebackgroundcolor=cGreen,frametitlefontcolor=white,
  innertopmargin=6pt,innerbottommargin=7pt,skipabove=9pt,skipbelow=9pt,
  innerleftmargin=9pt,innerrightmargin=9pt,roundcorner=3pt]{thmbox}   % משפט/טענה
\newmdenv[linecolor=cAmber,backgroundcolor=cAmberBg,linewidth=1.1pt,
  frametitlebackgroundcolor=cAmber,frametitlefontcolor=white,
  innertopmargin=6pt,innerbottommargin=7pt,skipabove=9pt,skipbelow=9pt,
  innerleftmargin=9pt,innerrightmargin=9pt,roundcorner=3pt]{notebox}  % הערה
\newmdenv[linecolor=cPurple,backgroundcolor=cPurpleBg,linewidth=1.1pt,
  frametitlebackgroundcolor=cPurple,frametitlefontcolor=white,
  innertopmargin=6pt,innerbottommargin=7pt,skipabove=9pt,skipbelow=9pt,
  innerleftmargin=9pt,innerrightmargin=9pt,roundcorner=3pt]{exbox}    % דוגמה
\newmdenv[linecolor=cRed,backgroundcolor=cRedBg,linewidth=1.1pt,
  frametitlebackgroundcolor=cRed,frametitlefontcolor=white,
  innertopmargin=6pt,innerbottommargin=7pt,skipabove=9pt,skipbelow=9pt,
  innerleftmargin=9pt,innerrightmargin=9pt,roundcorner=3pt]{warnbox}  % טעות נפוצה
\newmdenv[linecolor=cTeal,backgroundcolor=cTealBg,linewidth=1.1pt,
  frametitlebackgroundcolor=cTeal,frametitlefontcolor=white,
  innertopmargin=6pt,innerbottommargin=7pt,skipabove=9pt,skipbelow=9pt,
  innerleftmargin=9pt,innerrightmargin=9pt,roundcorner=3pt]{keybox}   % רעיון מרכזי

% ============================================================================
% Unified TABLE presentation: subtle uniform slate frame + auto-numbered
% caption rendered INSIDE the frame. Use:
%   \begin{tablebox}\centering\small  <tabular>  \tabcap{<description>}\end{tablebox}
% Replaces the buggy \caption* (which emitted a stray "טבלה N : *" line) and
% gives every table one identical frame + numbered caption "טבלה N.M.".
% Optional intro text may be placed at the TOP, before the tabular, inside the box.
% ============================================================================
\definecolor{cTblLine}{HTML}{B9C4D2}  \definecolor{cTblBg}{HTML}{FAFBFC}
% v3.23: class-aware — under `article` there is no chapter counter (the promised
% report->article switch used to die with "No counter 'chapter' defined").
\makeatletter
\@ifundefined{c@chapter}{%
  \newcounter{tbl}[section]\renewcommand{\thetbl}{\arabic{tbl}}%
}{%
  \newcounter{tbl}[chapter]\renewcommand{\thetbl}{\thechapter.\arabic{tbl}}%
}
\makeatother
\newmdenv[linecolor=cTblLine,backgroundcolor=cTblBg,linewidth=0.9pt,
  innertopmargin=9pt,innerbottommargin=8pt,skipabove=11pt,skipbelow=11pt,
  innerleftmargin=11pt,innerrightmargin=11pt,roundcorner=4pt]{tablebox}
\newcommand{\tabcap}[1]{\refstepcounter{tbl}\par\addvspace{6pt}%
  {\centering\small\textcolor{cGray}{\textbf{טבלה~\thetbl.}}\enspace #1\par}}

% ============================================================================
% LTR islands.  EVERY tikzpicture / quantikz MUST live inside one, or the whole
% figure's coordinates and Latin labels flip. (Hebrew node text still uses \h{}.)
% ============================================================================
\newenvironment{qcirc}{\par\medskip\begin{otherlanguage}{english}\begin{center}}
  {\end{center}\end{otherlanguage}\par\medskip}
\newenvironment{tikzpic}{\par\medskip\begin{otherlanguage}{english}\begin{center}}
  {\end{center}\end{otherlanguage}\par\medskip}

% ============================================================================
% Code listings.  ★★ v3.4 THE BOX-OVERFLOW FIX:
%   resetmargins=true   → margins measured from the ENCLOSING box, not the page
%                         (without it, a listing inside an mdframed box spills
%                          its frame + line-numbers out past the box border).
%   framexleftmargin=0  → the frame (left bar) does not extend left of the code.
%   xleftmargin keeps a small indent so the line numbers have room INSIDE.
% asmblock wraps the listing in an LTR island (code is LTR).
% ============================================================================
\lstdefinestyle{asm}{
  basicstyle=\ttfamily\small,
  backgroundcolor=\color{cBlueBg!55},
  frame=leftline,framerule=2.5pt,rulecolor=\color{cBlue},
  numbers=left,numberstyle=\tiny\color{cGray},numbersep=8pt,
  xleftmargin=2.2em,framexleftmargin=0pt,framexrightmargin=0pt,resetmargins=true,
  breaklines=true,columns=fullflexible,keepspaces=true,
  showstringspaces=false,aboveskip=8pt,belowskip=8pt,
  commentstyle=\color{cGreen},
}
\lstdefinestyle{cstyle}{
  language=C,basicstyle=\ttfamily\small,
  backgroundcolor=\color{cGreenBg!55},
  frame=leftline,framerule=2.5pt,rulecolor=\color{cGreen},
  numbers=left,numberstyle=\tiny\color{cGray},numbersep=8pt,
  xleftmargin=2.2em,framexleftmargin=0pt,framexrightmargin=0pt,resetmargins=true,
  breaklines=true,columns=fullflexible,keepspaces=true,
  showstringspaces=false,keywordstyle=\color{cBlue}\bfseries,
  commentstyle=\color{cGray}\itshape,aboveskip=8pt,belowskip=8pt,
}
\newenvironment{asmblock}{\par\begin{otherlanguage}{english}}{\end{otherlanguage}\par}

\usepackage[hidelinks,unicode]{hyperref}   % LAST, after colors. unicode = Hebrew bookmarks.

\setcounter{tocdepth}{1}
\setcounter{secnumdepth}{2}
     % (or paste these lines at the top)
%   \begin{document}  ...  \end{document}
%
% Compile with:  lualatex -interaction=nonstopmode doc.tex   (twice, for TOC)
% Requires the Hebrew fonts installed (run scripts/setup_fonts.sh first).
%
% SHORT doc (exam + solution): change `report` -> `article`.
% STUDY doc meant to be read/annotated: 12pt + \linespread{1.3} below are
%   deliberate readability choices (NOT padding) — keep or trim to taste.
% ============================================================================
\documentclass[12pt]{report}          % 12pt: readable for a study guide; 11pt is fine too
\usepackage[no-math]{fontspec}        % (1) ★ no-math — THE critical line (see recipe.md)
\usepackage[bidi=basic]{babel}        % (2) RTL handling built into the engine
\babelprovide[import,main]{hebrew}    % (3) Hebrew = main language, default RTL
\babelprovide[import]{english}        % (4) English secondary, for correct LTR runs
% Latin fallback registered BEFORE \babelfont so a lone English word in Hebrew
% prose does not tofu. Fonts resolved by FAMILY NAME (portable; no hard path).
\directlua{luaotfload.add_fallback("hebfb",{"DejaVuSerif:mode=node;"})}
% v3.22: Frank Ruhl Libre ships NO italic face. Instead of silently-upright
% \emph (v<=3.20) or bold-as-emphasis (v3.21), synthesize a slanted face with
% FakeSlant — real, VISIBLE italic emphasis with correct \emph nesting
% semantics (inner \emph toggles back upright). Verified empirically at
% 300dpi: slant clearly visible, 0 errors, 0 missing glyphs.
\babelfont{rm}[RawFeature={fallback=hebfb},
  ItalicFont={Frank Ruhl Libre}, ItalicFeatures={FakeSlant=0.22},
  BoldItalicFont={Frank Ruhl Libre Bold}, BoldItalicFeatures={FakeSlant=0.22}
  ]{Frank Ruhl Libre}
\babelfont{sf}[RawFeature={fallback=hebfb}]{Heebo}
\babelfont{tt}[RawFeature={fallback=hebfb}]{DejaVu Sans Mono}

\usepackage{amsmath,amssymb,mathtools}
\usepackage[margin=2.3cm,headheight=14pt]{geometry}
\linespread{1.3}                      % readable line spacing for study/annotation docs
\usepackage{xcolor}
\usepackage{mdframed}
\usepackage{listings}                 % code listings (assembly / C / pseudo-code)
\usepackage{tikz}
\usetikzlibrary{arrows.meta,positioning,calc,shapes.geometric,fit,backgrounds}
\usepackage{pgfplots}\pgfplotsset{compat=1.18}
\usepackage{quantikz}                 % quantum circuits (remove if unused)
\usepackage{booktabs}
\usepackage{longtable}
\usepackage{array}
\usepackage{float}                    % [H] = place table exactly here (reference doc: no drifting)
\usepackage{enumitem}
\usepackage{needspace}

% ============================================================================
% ★★ v3.4 — THE DIRECTION-ISOLATION TRIO (the heart of mixing Hebrew + Latin)
% Mental model:
%   \h{...}    Hebrew (RTL) island      — Hebrew text inside an LTR context
%                                          (TikZ nodes, the english `tikzpic`).
%   \en{...}   English (LTR) isolate    — Latin text / parentheses inside a
%                                          Hebrew (RTL) context. KEEPS PARENS
%                                          AND ORDER CORRECT.
%   \code{...} \en + monospace          — code inside Hebrew / inside a frametitle.
% CRITICAL: never wrap Latin in \h{} (it is for HEBREW). Inside a TikZ node,
% Latin in \h{} renders REVERSED (Datapath->htapataD, CPU->UPC) and parens
% mirror. Isolate every Latin/paren chunk with \en{} or \code{}. See
% references/figures-boxes-listings.md.
% ---- TEXT INSIDE MATH (v3.8) ----------------------------------------------
% Put ANY text that sits inside math in a \text{}. The ONLY thing that is
% invalid in math is a bare \h/\en (= \foreignlanguage) used DIRECTLY, outside
% \text -> "\rmfamily invalid in math mode". Inside \text:
%   Hebrew label           -> \text{...}            (bare Hebrew renders fine;
%                              \h is OPTIONAL/harmless: \text{\h{...}} also works
%                              and both render identically + scale in scripts)
%   plain Latin word        -> \text{...}            (no marker needed)
%   Latin WITH ( ) or [ ]   -> \text{\en{...}}       (\en forces LTR; a plain
%                              \text mirrors the parens in RTL -> "Spec (TNR(" )
%   f(Hebrew/text arg)      -> \text{f}(\text{...})  -- keep the parens at MATH
%                              level, NOT \text{f(arg)} (parens inside the box mirror)
%   code in math            -> \mathtt{...}  (or \texttt{}, math-guarded below)
% NEVER write \en{...} or \h{...} straight into $...$ / \[..\] (outside \text);
% always via \text{}. (\mbox{...} also works but does not scale in scripts.)
% ============================================================================
\newcommand{\h}[1]{\foreignlanguage{hebrew}{#1}}
\newcommand{\en}[1]{\foreignlanguage{english}{#1}}

% ============================================================================
% ★★ — fixes for defects that are INVISIBLE in the log and in the checks
% and only surface on a VISUAL RENDER of the PDF. (Each verified empirically.)
% ----------------------------------------------------------------------------
% v3.22: \emph is now REAL slanted emphasis (FakeSlant italic face defined on
% the \babelfont{rm} line above), replacing the v3.21 \emph->\textbf hack: it
% preserves \emph's toggle semantics and keeps bold free for its own role.
% Prefer bold-as-emphasis for your doc? one line: \renewcommand{\emph}[1]{\textbf{#1}}
% page number: force the digits LTR so they never reverse (18 -> 81) in the RTL
%   footer. That reversal is a NON-DETERMINISTIC bidi glitch (observed once in
%   ~80 pp); this LuaLaTeX build is not byte-reproducible, so the defensive
%   LTR wrap is the only reliable guard. Latin digits are always LTR anyway.
\renewcommand{\thepage}{\foreignlanguage{english}{\arabic{page}}}
% emergencystretch: Hebrew does not hyphenate, so an unbreakable LTR island
%   (\en{...} / an inline math run) landing at the margin can leave a small
%   (~2-13pt) overfull with no legal breakpoint. This rescues ONLY those
%   inline-prose tie-breaks; it is a ceiling used only on offending lines and
%   does NOT touch display-math / listing / table overflow (those still surface
%   in the log for the content-level overflow pass). Verified: baseline overfull
%   -> 0 with this set, with no change to pagination or underfull.
\setlength{\emergencystretch}{3em}
% em-dash: --- does NOT become — under Frank Ruhl (the font lacks tlig; verified
%   the width is identical with/without Ligatures=TeX, while the mechanism works
%   on Latin Modern). TYPE THE CHARACTER — (U+2014) DIRECTLY; likewise – (U+2013).
%   check_bidi_figures.py flags a --/--- that is adjacent to Hebrew.
% ============================================================================

% ---- STRUCTURAL DIRECTION-SAFETY for \texttt (v3.7) -------------------------
% A parenthesis, bracket, or $ INSIDE \texttt{} mirrors/migrates in RTL prose
% (MIPS 0($s0)->"(s0$)0", $rs->"rs$", arr[i]->"[arr[i") because \texttt sets the
% font but NOT the direction. Redefining \texttt to LTR-isolate its content fixes
% the WHOLE class at once (parens, brackets, $) so plain \texttt{lw $t0,0($s0)}
% is safe — no need to remember \code{}. Guarded by \ifmmode because
% \foreignlanguage is invalid in math (there, plain monospace is already LTR).
% NB: \texttt{Hebrew} stays tofu — the monospace font has no Hebrew anyway; put
% Hebrew in \h{}/normal text, never in \texttt{}. To disable, delete this block.
\let\origtexttt\texttt
\renewcommand{\texttt}[1]{\ifmmode\origtexttt{#1}\else\foreignlanguage{english}{\origtexttt{#1}}\fi}
\newcommand{\code}[1]{\ifmmode\origtexttt{#1}\else\foreignlanguage{english}{\origtexttt{#1}}\fi}

% ---- a thin decorative rule helper ----
\newcommand{\sep}{\par\vspace{2pt}\noindent\textcolor{cRule}{\rule{\linewidth}{0.6pt}}\par\vspace{4pt}}

% ---- colors ----
\definecolor{cBlue}{HTML}{1F4E79}   \definecolor{cBlueBg}{HTML}{EAF1F8}
\definecolor{cGreen}{HTML}{2E6B4F}  \definecolor{cGreenBg}{HTML}{EAF4EE}
\definecolor{cAmber}{HTML}{9C6A00}  \definecolor{cAmberBg}{HTML}{FBF3E2}
\definecolor{cPurple}{HTML}{5B2C6F} \definecolor{cPurpleBg}{HTML}{F3ECF7}
\definecolor{cTeal}{HTML}{0E6E6E}   \definecolor{cTealBg}{HTML}{E6F4F4}
\definecolor{cRed}{HTML}{8E2B2B}    \definecolor{cRedBg}{HTML}{FBEAEA}
\definecolor{cGray}{HTML}{555555}   \definecolor{cRule}{HTML}{1F4E79}

% ---- math operators ----
\DeclareMathOperator{\Tr}{Tr}
\DeclareMathOperator{\rank}{rank}
\DeclareMathOperator{\diag}{diag}
\DeclareMathOperator{\sgn}{sgn}

% ---- bra-ket (quantikz defines \ket/\bra/\braket; add the rest safely) ----
\providecommand{\ket}[1]{\left|#1\right\rangle}
\providecommand{\bra}[1]{\left\langle#1\right|}
\providecommand{\braket}[2]{\left\langle#1\middle|#2\right\rangle}
\newcommand{\dyad}[1]{\left|#1\right\rangle\!\left\langle#1\right|}
% v3.21: general outer product |a><b| and expectation value <A> — agents reach
% for these by reflex (esp. \expval); quantikz defines neither, so their absence
% threw "Undefined control sequence" in a real build. \providecommand = safe.
\providecommand{\ketbra}[2]{\left|#1\right\rangle\!\left\langle#2\right|}
\providecommand{\expval}[1]{\left\langle#1\right\rangle}
\newcommand{\abs}[1]{\left|#1\right|}
\newcommand{\norm}[1]{\left\lVert#1\right\rVert}
\newcommand{\half}{\tfrac12}
\newcommand{\R}{\mathbb{R}}\newcommand{\C}{\mathbb{C}}\newcommand{\Z}{\mathbb{Z}}
% identity operator. NB: \mathbb{1} silently renders a WRONG glyph (msbm has no
% blackboard digits). Use \mathbb{I} = 𝕀 (works with amssymb, no extra package).
% Do NOT "fix" this back to \mathbb{1}.
\newcommand{\id}{\mathbb{I}}

% ============================================================================
% Colored boxes — mdframed (NOT tcolorbox, which breaks under bidi).
% ★★ v3.4: innerleftmargin/innerrightmargin give code listings room so their
% frame/line-numbers do not spill out the box border (see listing styles below).
% ============================================================================
\newmdenv[linecolor=cBlue,backgroundcolor=cBlueBg,linewidth=1.1pt,
  frametitlebackgroundcolor=cBlue,frametitlefontcolor=white,
  innertopmargin=6pt,innerbottommargin=7pt,skipabove=9pt,skipbelow=9pt,
  innerleftmargin=9pt,innerrightmargin=9pt,roundcorner=3pt]{defbox}   % הגדרה
\newmdenv[linecolor=cGreen,backgroundcolor=cGreenBg,linewidth=1.1pt,
  frametitlebackgroundcolor=cGreen,frametitlefontcolor=white,
  innertopmargin=6pt,innerbottommargin=7pt,skipabove=9pt,skipbelow=9pt,
  innerleftmargin=9pt,innerrightmargin=9pt,roundcorner=3pt]{thmbox}   % משפט/טענה
\newmdenv[linecolor=cAmber,backgroundcolor=cAmberBg,linewidth=1.1pt,
  frametitlebackgroundcolor=cAmber,frametitlefontcolor=white,
  innertopmargin=6pt,innerbottommargin=7pt,skipabove=9pt,skipbelow=9pt,
  innerleftmargin=9pt,innerrightmargin=9pt,roundcorner=3pt]{notebox}  % הערה
\newmdenv[linecolor=cPurple,backgroundcolor=cPurpleBg,linewidth=1.1pt,
  frametitlebackgroundcolor=cPurple,frametitlefontcolor=white,
  innertopmargin=6pt,innerbottommargin=7pt,skipabove=9pt,skipbelow=9pt,
  innerleftmargin=9pt,innerrightmargin=9pt,roundcorner=3pt]{exbox}    % דוגמה
\newmdenv[linecolor=cRed,backgroundcolor=cRedBg,linewidth=1.1pt,
  frametitlebackgroundcolor=cRed,frametitlefontcolor=white,
  innertopmargin=6pt,innerbottommargin=7pt,skipabove=9pt,skipbelow=9pt,
  innerleftmargin=9pt,innerrightmargin=9pt,roundcorner=3pt]{warnbox}  % טעות נפוצה
\newmdenv[linecolor=cTeal,backgroundcolor=cTealBg,linewidth=1.1pt,
  frametitlebackgroundcolor=cTeal,frametitlefontcolor=white,
  innertopmargin=6pt,innerbottommargin=7pt,skipabove=9pt,skipbelow=9pt,
  innerleftmargin=9pt,innerrightmargin=9pt,roundcorner=3pt]{keybox}   % רעיון מרכזי

% ============================================================================
% Unified TABLE presentation: subtle uniform slate frame + auto-numbered
% caption rendered INSIDE the frame. Use:
%   \begin{tablebox}\centering\small  <tabular>  \tabcap{<description>}\end{tablebox}
% Replaces the buggy \caption* (which emitted a stray "טבלה N : *" line) and
% gives every table one identical frame + numbered caption "טבלה N.M.".
% Optional intro text may be placed at the TOP, before the tabular, inside the box.
% ============================================================================
\definecolor{cTblLine}{HTML}{B9C4D2}  \definecolor{cTblBg}{HTML}{FAFBFC}
% v3.23: class-aware — under `article` there is no chapter counter (the promised
% report->article switch used to die with "No counter 'chapter' defined").
\makeatletter
\@ifundefined{c@chapter}{%
  \newcounter{tbl}[section]\renewcommand{\thetbl}{\arabic{tbl}}%
}{%
  \newcounter{tbl}[chapter]\renewcommand{\thetbl}{\thechapter.\arabic{tbl}}%
}
\makeatother
\newmdenv[linecolor=cTblLine,backgroundcolor=cTblBg,linewidth=0.9pt,
  innertopmargin=9pt,innerbottommargin=8pt,skipabove=11pt,skipbelow=11pt,
  innerleftmargin=11pt,innerrightmargin=11pt,roundcorner=4pt]{tablebox}
\newcommand{\tabcap}[1]{\refstepcounter{tbl}\par\addvspace{6pt}%
  {\centering\small\textcolor{cGray}{\textbf{טבלה~\thetbl.}}\enspace #1\par}}

% ============================================================================
% LTR islands.  EVERY tikzpicture / quantikz MUST live inside one, or the whole
% figure's coordinates and Latin labels flip. (Hebrew node text still uses \h{}.)
% ============================================================================
\newenvironment{qcirc}{\par\medskip\begin{otherlanguage}{english}\begin{center}}
  {\end{center}\end{otherlanguage}\par\medskip}
\newenvironment{tikzpic}{\par\medskip\begin{otherlanguage}{english}\begin{center}}
  {\end{center}\end{otherlanguage}\par\medskip}

% ============================================================================
% Code listings.  ★★ v3.4 THE BOX-OVERFLOW FIX:
%   resetmargins=true   → margins measured from the ENCLOSING box, not the page
%                         (without it, a listing inside an mdframed box spills
%                          its frame + line-numbers out past the box border).
%   framexleftmargin=0  → the frame (left bar) does not extend left of the code.
%   xleftmargin keeps a small indent so the line numbers have room INSIDE.
% asmblock wraps the listing in an LTR island (code is LTR).
% ============================================================================
\lstdefinestyle{asm}{
  basicstyle=\ttfamily\small,
  backgroundcolor=\color{cBlueBg!55},
  frame=leftline,framerule=2.5pt,rulecolor=\color{cBlue},
  numbers=left,numberstyle=\tiny\color{cGray},numbersep=8pt,
  xleftmargin=2.2em,framexleftmargin=0pt,framexrightmargin=0pt,resetmargins=true,
  breaklines=true,columns=fullflexible,keepspaces=true,
  showstringspaces=false,aboveskip=8pt,belowskip=8pt,
  commentstyle=\color{cGreen},
}
\lstdefinestyle{cstyle}{
  language=C,basicstyle=\ttfamily\small,
  backgroundcolor=\color{cGreenBg!55},
  frame=leftline,framerule=2.5pt,rulecolor=\color{cGreen},
  numbers=left,numberstyle=\tiny\color{cGray},numbersep=8pt,
  xleftmargin=2.2em,framexleftmargin=0pt,framexrightmargin=0pt,resetmargins=true,
  breaklines=true,columns=fullflexible,keepspaces=true,
  showstringspaces=false,keywordstyle=\color{cBlue}\bfseries,
  commentstyle=\color{cGray}\itshape,aboveskip=8pt,belowskip=8pt,
}
\newenvironment{asmblock}{\par\begin{otherlanguage}{english}}{\end{otherlanguage}\par}

\usepackage[hidelinks,unicode]{hyperref}   % LAST, after colors. unicode = Hebrew bookmarks.

\setcounter{tocdepth}{1}
\setcounter{secnumdepth}{2}
     % (or paste these lines at the top)
%   \begin{document}  ...  \end{document}
%
% Compile with:  lualatex -interaction=nonstopmode doc.tex   (twice, for TOC)
% Requires the Hebrew fonts installed (run scripts/setup_fonts.sh first).
%
% SHORT doc (exam + solution): change `report` -> `article`.
% STUDY doc meant to be read/annotated: 12pt + \linespread{1.3} below are
%   deliberate readability choices (NOT padding) — keep or trim to taste.
% ============================================================================
\documentclass[12pt]{report}          % 12pt: readable for a study guide; 11pt is fine too
\usepackage[no-math]{fontspec}        % (1) ★ no-math — THE critical line (see recipe.md)
\usepackage[bidi=basic]{babel}        % (2) RTL handling built into the engine
\babelprovide[import,main]{hebrew}    % (3) Hebrew = main language, default RTL
\babelprovide[import]{english}        % (4) English secondary, for correct LTR runs
% Latin fallback registered BEFORE \babelfont so a lone English word in Hebrew
% prose does not tofu. Fonts resolved by FAMILY NAME (portable; no hard path).
\directlua{luaotfload.add_fallback("hebfb",{"DejaVuSerif:mode=node;"})}
% v3.22: Frank Ruhl Libre ships NO italic face. Instead of silently-upright
% \emph (v<=3.20) or bold-as-emphasis (v3.21), synthesize a slanted face with
% FakeSlant — real, VISIBLE italic emphasis with correct \emph nesting
% semantics (inner \emph toggles back upright). Verified empirically at
% 300dpi: slant clearly visible, 0 errors, 0 missing glyphs.
\babelfont{rm}[RawFeature={fallback=hebfb},
  ItalicFont={Frank Ruhl Libre}, ItalicFeatures={FakeSlant=0.22},
  BoldItalicFont={Frank Ruhl Libre Bold}, BoldItalicFeatures={FakeSlant=0.22}
  ]{Frank Ruhl Libre}
\babelfont{sf}[RawFeature={fallback=hebfb}]{Heebo}
\babelfont{tt}[RawFeature={fallback=hebfb}]{DejaVu Sans Mono}

\usepackage{amsmath,amssymb,mathtools}
\usepackage[margin=2.3cm,headheight=14pt]{geometry}
\linespread{1.3}                      % readable line spacing for study/annotation docs
\usepackage{xcolor}
\usepackage{mdframed}
\usepackage{listings}                 % code listings (assembly / C / pseudo-code)
\usepackage{tikz}
\usetikzlibrary{arrows.meta,positioning,calc,shapes.geometric,fit,backgrounds}
\usepackage{pgfplots}\pgfplotsset{compat=1.18}
\usepackage{quantikz}                 % quantum circuits (remove if unused)
\usepackage{booktabs}
\usepackage{longtable}
\usepackage{array}
\usepackage{float}                    % [H] = place table exactly here (reference doc: no drifting)
\usepackage{enumitem}
\usepackage{needspace}

% ============================================================================
% ★★ v3.4 — THE DIRECTION-ISOLATION TRIO (the heart of mixing Hebrew + Latin)
% Mental model:
%   \h{...}    Hebrew (RTL) island      — Hebrew text inside an LTR context
%                                          (TikZ nodes, the english `tikzpic`).
%   \en{...}   English (LTR) isolate    — Latin text / parentheses inside a
%                                          Hebrew (RTL) context. KEEPS PARENS
%                                          AND ORDER CORRECT.
%   \code{...} \en + monospace          — code inside Hebrew / inside a frametitle.
% CRITICAL: never wrap Latin in \h{} (it is for HEBREW). Inside a TikZ node,
% Latin in \h{} renders REVERSED (Datapath->htapataD, CPU->UPC) and parens
% mirror. Isolate every Latin/paren chunk with \en{} or \code{}. See
% references/figures-boxes-listings.md.
% ---- TEXT INSIDE MATH (v3.8) ----------------------------------------------
% Put ANY text that sits inside math in a \text{}. The ONLY thing that is
% invalid in math is a bare \h/\en (= \foreignlanguage) used DIRECTLY, outside
% \text -> "\rmfamily invalid in math mode". Inside \text:
%   Hebrew label           -> \text{...}            (bare Hebrew renders fine;
%                              \h is OPTIONAL/harmless: \text{\h{...}} also works
%                              and both render identically + scale in scripts)
%   plain Latin word        -> \text{...}            (no marker needed)
%   Latin WITH ( ) or [ ]   -> \text{\en{...}}       (\en forces LTR; a plain
%                              \text mirrors the parens in RTL -> "Spec (TNR(" )
%   f(Hebrew/text arg)      -> \text{f}(\text{...})  -- keep the parens at MATH
%                              level, NOT \text{f(arg)} (parens inside the box mirror)
%   code in math            -> \mathtt{...}  (or \texttt{}, math-guarded below)
% NEVER write \en{...} or \h{...} straight into $...$ / \[..\] (outside \text);
% always via \text{}. (\mbox{...} also works but does not scale in scripts.)
% ============================================================================
\newcommand{\h}[1]{\foreignlanguage{hebrew}{#1}}
\newcommand{\en}[1]{\foreignlanguage{english}{#1}}

% ============================================================================
% ★★ — fixes for defects that are INVISIBLE in the log and in the checks
% and only surface on a VISUAL RENDER of the PDF. (Each verified empirically.)
% ----------------------------------------------------------------------------
% v3.22: \emph is now REAL slanted emphasis (FakeSlant italic face defined on
% the \babelfont{rm} line above), replacing the v3.21 \emph->\textbf hack: it
% preserves \emph's toggle semantics and keeps bold free for its own role.
% Prefer bold-as-emphasis for your doc? one line: \renewcommand{\emph}[1]{\textbf{#1}}
% page number: force the digits LTR so they never reverse (18 -> 81) in the RTL
%   footer. That reversal is a NON-DETERMINISTIC bidi glitch (observed once in
%   ~80 pp); this LuaLaTeX build is not byte-reproducible, so the defensive
%   LTR wrap is the only reliable guard. Latin digits are always LTR anyway.
\renewcommand{\thepage}{\foreignlanguage{english}{\arabic{page}}}
% emergencystretch: Hebrew does not hyphenate, so an unbreakable LTR island
%   (\en{...} / an inline math run) landing at the margin can leave a small
%   (~2-13pt) overfull with no legal breakpoint. This rescues ONLY those
%   inline-prose tie-breaks; it is a ceiling used only on offending lines and
%   does NOT touch display-math / listing / table overflow (those still surface
%   in the log for the content-level overflow pass). Verified: baseline overfull
%   -> 0 with this set, with no change to pagination or underfull.
\setlength{\emergencystretch}{3em}
% em-dash: --- does NOT become — under Frank Ruhl (the font lacks tlig; verified
%   the width is identical with/without Ligatures=TeX, while the mechanism works
%   on Latin Modern). TYPE THE CHARACTER — (U+2014) DIRECTLY; likewise – (U+2013).
%   check_bidi_figures.py flags a --/--- that is adjacent to Hebrew.
% ============================================================================

% ---- STRUCTURAL DIRECTION-SAFETY for \texttt (v3.7) -------------------------
% A parenthesis, bracket, or $ INSIDE \texttt{} mirrors/migrates in RTL prose
% (MIPS 0($s0)->"(s0$)0", $rs->"rs$", arr[i]->"[arr[i") because \texttt sets the
% font but NOT the direction. Redefining \texttt to LTR-isolate its content fixes
% the WHOLE class at once (parens, brackets, $) so plain \texttt{lw $t0,0($s0)}
% is safe — no need to remember \code{}. Guarded by \ifmmode because
% \foreignlanguage is invalid in math (there, plain monospace is already LTR).
% NB: \texttt{Hebrew} stays tofu — the monospace font has no Hebrew anyway; put
% Hebrew in \h{}/normal text, never in \texttt{}. To disable, delete this block.
\let\origtexttt\texttt
\renewcommand{\texttt}[1]{\ifmmode\origtexttt{#1}\else\foreignlanguage{english}{\origtexttt{#1}}\fi}
\newcommand{\code}[1]{\ifmmode\origtexttt{#1}\else\foreignlanguage{english}{\origtexttt{#1}}\fi}

% ---- a thin decorative rule helper ----
\newcommand{\sep}{\par\vspace{2pt}\noindent\textcolor{cRule}{\rule{\linewidth}{0.6pt}}\par\vspace{4pt}}

% ---- colors ----
\definecolor{cBlue}{HTML}{1F4E79}   \definecolor{cBlueBg}{HTML}{EAF1F8}
\definecolor{cGreen}{HTML}{2E6B4F}  \definecolor{cGreenBg}{HTML}{EAF4EE}
\definecolor{cAmber}{HTML}{9C6A00}  \definecolor{cAmberBg}{HTML}{FBF3E2}
\definecolor{cPurple}{HTML}{5B2C6F} \definecolor{cPurpleBg}{HTML}{F3ECF7}
\definecolor{cTeal}{HTML}{0E6E6E}   \definecolor{cTealBg}{HTML}{E6F4F4}
\definecolor{cRed}{HTML}{8E2B2B}    \definecolor{cRedBg}{HTML}{FBEAEA}
\definecolor{cGray}{HTML}{555555}   \definecolor{cRule}{HTML}{1F4E79}

% ---- math operators ----
\DeclareMathOperator{\Tr}{Tr}
\DeclareMathOperator{\rank}{rank}
\DeclareMathOperator{\diag}{diag}
\DeclareMathOperator{\sgn}{sgn}

% ---- bra-ket (quantikz defines \ket/\bra/\braket; add the rest safely) ----
\providecommand{\ket}[1]{\left|#1\right\rangle}
\providecommand{\bra}[1]{\left\langle#1\right|}
\providecommand{\braket}[2]{\left\langle#1\middle|#2\right\rangle}
\newcommand{\dyad}[1]{\left|#1\right\rangle\!\left\langle#1\right|}
% v3.21: general outer product |a><b| and expectation value <A> — agents reach
% for these by reflex (esp. \expval); quantikz defines neither, so their absence
% threw "Undefined control sequence" in a real build. \providecommand = safe.
\providecommand{\ketbra}[2]{\left|#1\right\rangle\!\left\langle#2\right|}
\providecommand{\expval}[1]{\left\langle#1\right\rangle}
\newcommand{\abs}[1]{\left|#1\right|}
\newcommand{\norm}[1]{\left\lVert#1\right\rVert}
\newcommand{\half}{\tfrac12}
\newcommand{\R}{\mathbb{R}}\newcommand{\C}{\mathbb{C}}\newcommand{\Z}{\mathbb{Z}}
% identity operator. NB: \mathbb{1} silently renders a WRONG glyph (msbm has no
% blackboard digits). Use \mathbb{I} = 𝕀 (works with amssymb, no extra package).
% Do NOT "fix" this back to \mathbb{1}.
\newcommand{\id}{\mathbb{I}}

% ============================================================================
% Colored boxes — mdframed (NOT tcolorbox, which breaks under bidi).
% ★★ v3.4: innerleftmargin/innerrightmargin give code listings room so their
% frame/line-numbers do not spill out the box border (see listing styles below).
% ============================================================================
\newmdenv[linecolor=cBlue,backgroundcolor=cBlueBg,linewidth=1.1pt,
  frametitlebackgroundcolor=cBlue,frametitlefontcolor=white,
  innertopmargin=6pt,innerbottommargin=7pt,skipabove=9pt,skipbelow=9pt,
  innerleftmargin=9pt,innerrightmargin=9pt,roundcorner=3pt]{defbox}   % הגדרה
\newmdenv[linecolor=cGreen,backgroundcolor=cGreenBg,linewidth=1.1pt,
  frametitlebackgroundcolor=cGreen,frametitlefontcolor=white,
  innertopmargin=6pt,innerbottommargin=7pt,skipabove=9pt,skipbelow=9pt,
  innerleftmargin=9pt,innerrightmargin=9pt,roundcorner=3pt]{thmbox}   % משפט/טענה
\newmdenv[linecolor=cAmber,backgroundcolor=cAmberBg,linewidth=1.1pt,
  frametitlebackgroundcolor=cAmber,frametitlefontcolor=white,
  innertopmargin=6pt,innerbottommargin=7pt,skipabove=9pt,skipbelow=9pt,
  innerleftmargin=9pt,innerrightmargin=9pt,roundcorner=3pt]{notebox}  % הערה
\newmdenv[linecolor=cPurple,backgroundcolor=cPurpleBg,linewidth=1.1pt,
  frametitlebackgroundcolor=cPurple,frametitlefontcolor=white,
  innertopmargin=6pt,innerbottommargin=7pt,skipabove=9pt,skipbelow=9pt,
  innerleftmargin=9pt,innerrightmargin=9pt,roundcorner=3pt]{exbox}    % דוגמה
\newmdenv[linecolor=cRed,backgroundcolor=cRedBg,linewidth=1.1pt,
  frametitlebackgroundcolor=cRed,frametitlefontcolor=white,
  innertopmargin=6pt,innerbottommargin=7pt,skipabove=9pt,skipbelow=9pt,
  innerleftmargin=9pt,innerrightmargin=9pt,roundcorner=3pt]{warnbox}  % טעות נפוצה
\newmdenv[linecolor=cTeal,backgroundcolor=cTealBg,linewidth=1.1pt,
  frametitlebackgroundcolor=cTeal,frametitlefontcolor=white,
  innertopmargin=6pt,innerbottommargin=7pt,skipabove=9pt,skipbelow=9pt,
  innerleftmargin=9pt,innerrightmargin=9pt,roundcorner=3pt]{keybox}   % רעיון מרכזי

% ============================================================================
% Unified TABLE presentation: subtle uniform slate frame + auto-numbered
% caption rendered INSIDE the frame. Use:
%   \begin{tablebox}\centering\small  <tabular>  \tabcap{<description>}\end{tablebox}
% Replaces the buggy \caption* (which emitted a stray "טבלה N : *" line) and
% gives every table one identical frame + numbered caption "טבלה N.M.".
% Optional intro text may be placed at the TOP, before the tabular, inside the box.
% ============================================================================
\definecolor{cTblLine}{HTML}{B9C4D2}  \definecolor{cTblBg}{HTML}{FAFBFC}
% v3.23: class-aware — under `article` there is no chapter counter (the promised
% report->article switch used to die with "No counter 'chapter' defined").
\makeatletter
\@ifundefined{c@chapter}{%
  \newcounter{tbl}[section]\renewcommand{\thetbl}{\arabic{tbl}}%
}{%
  \newcounter{tbl}[chapter]\renewcommand{\thetbl}{\thechapter.\arabic{tbl}}%
}
\makeatother
\newmdenv[linecolor=cTblLine,backgroundcolor=cTblBg,linewidth=0.9pt,
  innertopmargin=9pt,innerbottommargin=8pt,skipabove=11pt,skipbelow=11pt,
  innerleftmargin=11pt,innerrightmargin=11pt,roundcorner=4pt]{tablebox}
\newcommand{\tabcap}[1]{\refstepcounter{tbl}\par\addvspace{6pt}%
  {\centering\small\textcolor{cGray}{\textbf{טבלה~\thetbl.}}\enspace #1\par}}

% ============================================================================
% LTR islands.  EVERY tikzpicture / quantikz MUST live inside one, or the whole
% figure's coordinates and Latin labels flip. (Hebrew node text still uses \h{}.)
% ============================================================================
\newenvironment{qcirc}{\par\medskip\begin{otherlanguage}{english}\begin{center}}
  {\end{center}\end{otherlanguage}\par\medskip}
\newenvironment{tikzpic}{\par\medskip\begin{otherlanguage}{english}\begin{center}}
  {\end{center}\end{otherlanguage}\par\medskip}

% ============================================================================
% Code listings.  ★★ v3.4 THE BOX-OVERFLOW FIX:
%   resetmargins=true   → margins measured from the ENCLOSING box, not the page
%                         (without it, a listing inside an mdframed box spills
%                          its frame + line-numbers out past the box border).
%   framexleftmargin=0  → the frame (left bar) does not extend left of the code.
%   xleftmargin keeps a small indent so the line numbers have room INSIDE.
% asmblock wraps the listing in an LTR island (code is LTR).
% ============================================================================
\lstdefinestyle{asm}{
  basicstyle=\ttfamily\small,
  backgroundcolor=\color{cBlueBg!55},
  frame=leftline,framerule=2.5pt,rulecolor=\color{cBlue},
  numbers=left,numberstyle=\tiny\color{cGray},numbersep=8pt,
  xleftmargin=2.2em,framexleftmargin=0pt,framexrightmargin=0pt,resetmargins=true,
  breaklines=true,columns=fullflexible,keepspaces=true,
  showstringspaces=false,aboveskip=8pt,belowskip=8pt,
  commentstyle=\color{cGreen},
}
\lstdefinestyle{cstyle}{
  language=C,basicstyle=\ttfamily\small,
  backgroundcolor=\color{cGreenBg!55},
  frame=leftline,framerule=2.5pt,rulecolor=\color{cGreen},
  numbers=left,numberstyle=\tiny\color{cGray},numbersep=8pt,
  xleftmargin=2.2em,framexleftmargin=0pt,framexrightmargin=0pt,resetmargins=true,
  breaklines=true,columns=fullflexible,keepspaces=true,
  showstringspaces=false,keywordstyle=\color{cBlue}\bfseries,
  commentstyle=\color{cGray}\itshape,aboveskip=8pt,belowskip=8pt,
}
\newenvironment{asmblock}{\par\begin{otherlanguage}{english}}{\end{otherlanguage}\par}

\usepackage[hidelinks,unicode]{hyperref}   % LAST, after colors. unicode = Hebrew bookmarks.

\setcounter{tocdepth}{1}
\setcounter{secnumdepth}{2}

\usepackage{graphicx}

\begin{document}
\begin{titlepage}
\centering\vspace*{3.8cm}
{\Huge\bfseries לקט מערכות וחישוב\par}
\vspace{1.0cm}
{\Large ארכיטקטורה / מערכות / אלגוריתמים / למידה / קוונטים / אבטחה\par}
\vspace{0.8cm}
{\large ראווה ל־\en{hebrew-lualatex-pdf} — לפי \texttt{content-style.md}\par}
\vspace{0.4cm}
{\large יעד: פרוזה + משוואות ממוספרות כשדרה; $\sim$1--2 קופסאות לפרק\par}
\vfill
\end{titlepage}

\tableofcontents
\clearpage

\part{ארכיטקטורת מעבדים}

\chapter{צנרת בסיסית}
\label{ch:pipeline}

היחידה נסגרת סביב \en{CPI}: הגדרה, מדד, ומה נשבר כשמתעלמים ממנו.

\begin{defbox}[frametitle={הגדרה — צנרת בסיסית}]
צנרת מפרקת הוראה לשלבים חופפים בזמן כך שבמצב יציב יוצאת הוראה בכל מחזור.
\end{defbox}

\begin{equation}
\label{eq:pipeline}
\mathrm{CPI}_{\mathrm{ideal}}=1
\end{equation}

המשמעות של~\eqref{eq:pipeline}: מודדים את \en{CPI} קודם, ורק אחר כך צוללים לרכיב.

\begin{tablebox}\centering\small
\begin{tabular}{@{}lll@{}}
\toprule
\textbf{מצב} & \textbf{סדר גודל} & \textbf{פעולה}\\
\midrule
תקין & $\sim 1\times$ & שמור מדידה\\
חשוד & $1.5$–$2\times$ & בודקים תלות\\
שבור & $\ge 2\times$ & תיקון מבני\\
\bottomrule
\end{tabular}
\tabcap{כיול מהיר ל־\en{CPI}.}
\end{tablebox}

רצף שממחיש לחץ על המדד:
\begin{asmblock}
\begin{lstlisting}[style=asm]
lw   $t0, 0($s0)
add  $t1, $t0, $s1   # depends on t0
\end{lstlisting}
\end{asmblock}
התלות על \code{\$t0} היא עלות ב־\en{CPI}, לא ״באג קוד״.

\begin{tikzpic}
\begin{tikzpicture}[font=\small,node distance=8mm,>={Latex[length=2mm]},
  box/.style={draw,thick,rounded corners,align=center,minimum width=2.4cm,minimum height=1.0cm,fill=cBlueBg}]
  \node[box] (a) {\en{In}};
  \node[box,right=of a] (b) {\en{CPI}};
  \node[box,right=of b,fill=cGreenBg] (c) {\en{Out}};
  \draw[->,thick] (a) -- (b);
  \draw[->,thick] (b) -- (c);
\end{tikzpicture}
\end{tikzpic}
\begin{center}\small איור — זרימה ל־\en{CPI}\end{center}

סדר עבודה:
\begin{enumerate}
\item[(\textbf{א})] מדוד את \en{CPI} לפי~\eqref{eq:pipeline}.
\item[(\textbf{ב})] זהה גורם שולט אחד.
\item[(\textbf{ג})] שנה אותו בלבד, ומדוד שוב.
\end{enumerate}

\chapter{סכסוכי נתונים}
\label{ch:hazards}

היחידה נסגרת סביב \en{RAW}: הגדרה, מדד, ומה נשבר כשמתעלמים ממנו.

\begin{warnbox}[frametitle={מלכודת — סכסוכי נתונים}]
גם עם קידום מלא, \code{lw} ואחריו שימוש דורשים \en{stall} של מחזור.
\end{warnbox}

\begin{equation}
\label{eq:hazards}
T_{\mathrm{stall}}=N_{\mathrm{load\text{-}use}}
\end{equation}

המשמעות של~\eqref{eq:hazards}: מודדים את \en{RAW} קודם, ורק אחר כך צוללים לרכיב.

\begin{equation}
\label{eq:hazards-check}
\Delta M \propto \Delta x_{\mathrm{dom}}
\end{equation}
כאן $M$ הוא \en{RAW};~\eqref{eq:hazards-check} מזכירה לחפש רגישות, לא מיקרו־אופטימיזציה.

סדר עבודה:
\begin{enumerate}
\item[(\textbf{א})] מדוד את \en{RAW} לפי~\eqref{eq:hazards}.
\item[(\textbf{ב})] זהה גורם שולט אחד.
\item[(\textbf{ג})] שנה אותו בלבד, ומדוד שוב.
\end{enumerate}

\chapter{סכסוכי בקרה}
\label{ch:control}

היחידה נסגרת סביב \en{branch}: הגדרה, מדד, ומה נשבר כשמתעלמים ממנו.

\begin{keybox}[frametitle={רעיון מפתח — סכסוכי בקרה}]
מחיר קפיצה שגויה הוא מכפלת ההסתברות בעונש במחזורים.
\end{keybox}

\begin{equation}
\label{eq:control}
P_{\mathrm{mispredict}}\cdot T_{\mathrm{penalty}}
\end{equation}

המשמעות של~\eqref{eq:control}: מודדים את \en{branch} קודם, ורק אחר כך צוללים לרכיב.

סדר עבודה:
\begin{enumerate}
\item[(\textbf{א})] מדוד את \en{branch} לפי~\eqref{eq:control}.
\item[(\textbf{ב})] זהה גורם שולט אחד.
\item[(\textbf{ג})] שנה אותו בלבד, ומדוד שוב.
\end{enumerate}

\chapter{קידום אוגרים}
\label{ch:forward}

היחידה נסגרת סביב \en{bypass}: הגדרה, מדד, ומה נשבר כשמתעלמים ממנו.

\begin{defbox}[frametitle={הגדרה — קידום אוגרים}]
קידום מעביר תוצאה מ־\en{EX/MEM} או \en{MEM/WB} בלי לחכות ל־\en{WB}.
\end{defbox}

\begin{equation}
\label{eq:forward}
T_{\mathrm{fwd}} \ll T_{\mathrm{stall}}
\end{equation}

המשמעות של~\eqref{eq:forward}: מודדים את \en{bypass} קודם, ורק אחר כך צוללים לרכיב.

\begin{equation}
\label{eq:forward-check}
\Delta M \propto \Delta x_{\mathrm{dom}}
\end{equation}
כאן $M$ הוא \en{bypass};~\eqref{eq:forward-check} מזכירה לחפש רגישות, לא מיקרו־אופטימיזציה.

סדר עבודה:
\begin{enumerate}
\item[(\textbf{א})] מדוד את \en{bypass} לפי~\eqref{eq:forward}.
\item[(\textbf{ב})] זהה גורם שולט אחד.
\item[(\textbf{ג})] שנה אותו בלבד, ומדוד שוב.
\end{enumerate}

\chapter{חיזוי קפיצות}
\label{ch:branchpred}

היחידה נסגרת סביב \en{BHT}: הגדרה, מדד, ומה נשבר כשמתעלמים ממנו.

\begin{notebox}[frametitle={הערה — חיזוי קפיצות}]
דיוק החיזוי $A$ הוא המשלים של שיעור השגיאה.
\end{notebox}

\begin{equation}
\label{eq:branchpred}
A = 1 - P_{\mathrm{mis}}
\end{equation}

המשמעות של~\eqref{eq:branchpred}: מודדים את \en{BHT} קודם, ורק אחר כך צוללים לרכיב.

\begin{tablebox}\centering\small
\begin{tabular}{@{}lll@{}}
\toprule
\textbf{מצב} & \textbf{סדר גודל} & \textbf{פעולה}\\
\midrule
תקין & $\sim 1\times$ & שמור מדידה\\
חשוד & $1.5$–$2\times$ & בודקים תלות\\
שבור & $\ge 2\times$ & תיקון מבני\\
\bottomrule
\end{tabular}
\tabcap{כיול מהיר ל־\en{BHT}.}
\end{tablebox}

סדר עבודה:
\begin{enumerate}
\item[(\textbf{א})] מדוד את \en{BHT} לפי~\eqref{eq:branchpred}.
\item[(\textbf{ב})] זהה גורם שולט אחד.
\item[(\textbf{ג})] שנה אותו בלבד, ומדוד שוב.
\end{enumerate}

\chapter{סופרסקלר}
\label{ch:superscalar}

היחידה נסגרת סביב \en{IPC}: הגדרה, מדד, ומה נשבר כשמתעלמים ממנו.

\begin{defbox}[frametitle={הגדרה — סופרסקלר}]
רוחב השיגור $W$ הוא חסם עליון ל־\en{IPC}, לא הערך הטיפוסי.
\end{defbox}

\begin{equation}
\label{eq:superscalar}
\mathrm{IPC}_{\mathrm{peak}} = W
\end{equation}

המשמעות של~\eqref{eq:superscalar}: מודדים את \en{IPC} קודם, ורק אחר כך צוללים לרכיב.

\begin{equation}
\label{eq:superscalar-check}
\Delta M \propto \Delta x_{\mathrm{dom}}
\end{equation}
כאן $M$ הוא \en{IPC};~\eqref{eq:superscalar-check} מזכירה לחפש רגישות, לא מיקרו־אופטימיזציה.

רצף שממחיש לחץ על המדד:
\begin{asmblock}
\begin{lstlisting}[style=asm]
lw   $t0, 0($s0)
add  $t1, $t0, $s1   # depends on t0
\end{lstlisting}
\end{asmblock}
התלות על \code{\$t0} היא עלות ב־\en{IPC}, לא ״באג קוד״.

סדר עבודה:
\begin{enumerate}
\item[(\textbf{א})] מדוד את \en{IPC} לפי~\eqref{eq:superscalar}.
\item[(\textbf{ב})] זהה גורם שולט אחד.
\item[(\textbf{ג})] שנה אותו בלבד, ומדוד שוב.
\end{enumerate}

\chapter{ביצוע מחוץ לסדר}
\label{ch:ooo}

היחידה נסגרת סביב \en{ROB}: הגדרה, מדד, ומה נשבר כשמתעלמים ממנו.

\begin{keybox}[frametitle={רעיון מפתח — ביצוע מחוץ לסדר}]
גודל ה־\en{ROB} חייב לכסות את חלון הביצוע $W\cdot L$.
\end{keybox}

\begin{equation}
\label{eq:ooo}
N_{\mathrm{ROB}} \ge W\cdot L
\end{equation}

המשמעות של~\eqref{eq:ooo}: מודדים את \en{ROB} קודם, ורק אחר כך צוללים לרכיב.

\begin{tikzpic}
\begin{tikzpicture}[font=\small,node distance=8mm,>={Latex[length=2mm]},
  box/.style={draw,thick,rounded corners,align=center,minimum width=2.4cm,minimum height=1.0cm,fill=cBlueBg}]
  \node[box] (a) {\en{In}};
  \node[box,right=of a] (b) {\en{ROB}};
  \node[box,right=of b,fill=cGreenBg] (c) {\en{Out}};
  \draw[->,thick] (a) -- (b);
  \draw[->,thick] (b) -- (c);
\end{tikzpicture}
\end{tikzpic}
\begin{center}\small איור — זרימה ל־\en{ROB}\end{center}

סדר עבודה:
\begin{enumerate}
\item[(\textbf{א})] מדוד את \en{ROB} לפי~\eqref{eq:ooo}.
\item[(\textbf{ב})] זהה גורם שולט אחד.
\item[(\textbf{ג})] שנה אותו בלבד, ומדוד שוב.
\end{enumerate}

\chapter{רישום מחדש}
\label{ch:rename}

היחידה נסגרת סביב \en{RAT}: הגדרה, מדד, ומה נשבר כשמתעלמים ממנו.

\begin{defbox}[frametitle={הגדרה — רישום מחדש}]
ה־\en{RAT} ממפה אוגר ארכיטקטוני לאוגר פיזיקלי ושובר תלויות־שם.
\end{defbox}

\begin{equation}
\label{eq:rename}
\mathrm{RAT}: a_r \mapsto p_i
\end{equation}

המשמעות של~\eqref{eq:rename}: מודדים את \en{RAT} קודם, ורק אחר כך צוללים לרכיב.

\begin{equation}
\label{eq:rename-check}
\Delta M \propto \Delta x_{\mathrm{dom}}
\end{equation}
כאן $M$ הוא \en{RAT};~\eqref{eq:rename-check} מזכירה לחפש רגישות, לא מיקרו־אופטימיזציה.

סדר עבודה:
\begin{enumerate}
\item[(\textbf{א})] מדוד את \en{RAT} לפי~\eqref{eq:rename}.
\item[(\textbf{ב})] זהה גורם שולט אחד.
\item[(\textbf{ג})] שנה אותו בלבד, ומדוד שוב.
\end{enumerate}


\chapter{חיבור — ארכיטקטורת מעבדים}
\label{ch:bridge-pipeline}

הפרקים בחלק ״ארכיטקטורת מעבדים״ חולקים משמעת אחת: מדד, נוסחה, מדידה.
הנוסחאות \\eqref{eq:pipeline}, \\eqref{eq:hazards}, \\eqref{eq:control}, \\eqref{eq:forward} אינן אוסף עובדות — הן חוזה מדידה.

\begin{keybox}[frametitle={חוזה החלק}]
לכל נוסחה ממוספרת בחלק כתוב ניסוי אחד שמאמת אותה. בלי ניסוי — היא דקורציה.
\end{keybox}

כששני מדדים מתנגשים, בוחרים לפי מטרת המערכת (השהיה מול מעבר־נתונים, דיוק מול עלות),
לא לפי הנוסחה היפה יותר.

\part{היררכיית זיכרון}

\chapter{מטמון L1}
\label{ch:l1}

היחידה נסגרת סביב \en{AMAT}: הגדרה, מדד, ומה נשבר כשמתעלמים ממנו.

\begin{keybox}[frametitle={רעיון מפתח — מטמון L1}]
\en{AMAT} הוא המדד היחיד שמאחד \en{hit time} ו־\en{miss rate}.
\end{keybox}

\begin{equation}
\label{eq:l1}
T_{\mathrm{avg}}=T_h+r_m T_m
\end{equation}

המשמעות של~\eqref{eq:l1}: מודדים את \en{AMAT} קודם, ורק אחר כך צוללים לרכיב.

\begin{tablebox}\centering\small
\begin{tabular}{@{}lll@{}}
\toprule
\textbf{מצב} & \textbf{סדר גודל} & \textbf{פעולה}\\
\midrule
תקין & $\sim 1\times$ & שמור מדידה\\
חשוד & $1.5$–$2\times$ & בודקים תלות\\
שבור & $\ge 2\times$ & תיקון מבני\\
\bottomrule
\end{tabular}
\tabcap{כיול מהיר ל־\en{AMAT}.}
\end{tablebox}

סדר עבודה:
\begin{enumerate}
\item[(\textbf{א})] מדוד את \en{AMAT} לפי~\eqref{eq:l1}.
\item[(\textbf{ב})] זהה גורם שולט אחד.
\item[(\textbf{ג})] שנה אותו בלבד, ומדוד שוב.
\end{enumerate}

\chapter{מטמון L2/L3}
\label{ch:l2}

היחידה נסגרת סביב \en{inclusive}: הגדרה, מדד, ומה נשבר כשמתעלמים ממנו.

\begin{defbox}[frametitle={הגדרה — מטמון L2/L3}]
נוסחת הרמות משרשרת עונשים לפי שיעורי החמצה בכל דרג.
\end{defbox}

\begin{equation}
\label{eq:l2}
T_{\mathrm{avg}}=T_{L1}+r_1(T_{L2}+r_2 T_{\mathrm{mem}})
\end{equation}

המשמעות של~\eqref{eq:l2}: מודדים את \en{inclusive} קודם, ורק אחר כך צוללים לרכיב.

\begin{equation}
\label{eq:l2-check}
\Delta M \propto \Delta x_{\mathrm{dom}}
\end{equation}
כאן $M$ הוא \en{inclusive};~\eqref{eq:l2-check} מזכירה לחפש רגישות, לא מיקרו־אופטימיזציה.

סדר עבודה:
\begin{enumerate}
\item[(\textbf{א})] מדוד את \en{inclusive} לפי~\eqref{eq:l2}.
\item[(\textbf{ב})] זהה גורם שולט אחד.
\item[(\textbf{ג})] שנה אותו בלבד, ומדוד שוב.
\end{enumerate}

\chapter{מדיניות החלפה}
\label{ch:repl}

היחידה נסגרת סביב \en{LRU}: הגדרה, מדד, ומה נשבר כשמתעלמים ממנו.

\begin{notebox}[frametitle={הערה — מדיניות החלפה}]
שלוש ה־C: compulsory, capacity, conflict — כל אחת דורשת טיפול אחר.
\end{notebox}

\begin{equation}
\label{eq:repl}
C_{\mathrm{miss}} = C_{\mathrm{comp}}+C_{\mathrm{cap}}+C_{\mathrm{conf}}
\end{equation}

המשמעות של~\eqref{eq:repl}: מודדים את \en{LRU} קודם, ורק אחר כך צוללים לרכיב.

רצף שממחיש לחץ על המדד:
\begin{asmblock}
\begin{lstlisting}[style=asm]
lw   $t0, 0($s0)
add  $t1, $t0, $s1   # depends on t0
\end{lstlisting}
\end{asmblock}
התלות על \code{\$t0} היא עלות ב־\en{LRU}, לא ״באג קוד״.

סדר עבודה:
\begin{enumerate}
\item[(\textbf{א})] מדוד את \en{LRU} לפי~\eqref{eq:repl}.
\item[(\textbf{ב})] זהה גורם שולט אחד.
\item[(\textbf{ג})] שנה אותו בלבד, ומדוד שוב.
\end{enumerate}

\chapter{מקומיות}
\label{ch:locality}

היחידה נסגרת סביב \en{temporal}: הגדרה, מדד, ומה נשבר כשמתעלמים ממנו.

\begin{defbox}[frametitle={הגדרה — מקומיות}]
מקומיות זמנית דועכת עם הזמן; מטמון מנצל את הזנב הקצר.
\end{defbox}

\begin{equation}
\label{eq:locality}
P(\mathrm{reuse}\mid t) \propto e^{-t/\tau}
\end{equation}

המשמעות של~\eqref{eq:locality}: מודדים את \en{temporal} קודם, ורק אחר כך צוללים לרכיב.

\begin{equation}
\label{eq:locality-check}
\Delta M \propto \Delta x_{\mathrm{dom}}
\end{equation}
כאן $M$ הוא \en{temporal};~\eqref{eq:locality-check} מזכירה לחפש רגישות, לא מיקרו־אופטימיזציה.

סדר עבודה:
\begin{enumerate}
\item[(\textbf{א})] מדוד את \en{temporal} לפי~\eqref{eq:locality}.
\item[(\textbf{ב})] זהה גורם שולט אחד.
\item[(\textbf{ג})] שנה אותו בלבד, ומדוד שוב.
\end{enumerate}

\chapter{TLB ותרגום}
\label{ch:tlb}

היחידה נסגרת סביב \en{VPN}: הגדרה, מדד, ומה נשבר כשמתעלמים ממנו.

\begin{keybox}[frametitle={רעיון מפתח — TLB ותרגום}]
החמצת \en{TLB} יקרה כי הולכים ל־\en{page walk}.
\end{keybox}

\begin{equation}
\label{eq:tlb}
T_{\mathrm{addr}}=T_{\mathrm{TLB}}+r_{\mathrm{TLB}}T_{\mathrm{PTW}}
\end{equation}

המשמעות של~\eqref{eq:tlb}: מודדים את \en{VPN} קודם, ורק אחר כך צוללים לרכיב.

\begin{tablebox}\centering\small
\begin{tabular}{@{}lll@{}}
\toprule
\textbf{מצב} & \textbf{סדר גודל} & \textbf{פעולה}\\
\midrule
תקין & $\sim 1\times$ & שמור מדידה\\
חשוד & $1.5$–$2\times$ & בודקים תלות\\
שבור & $\ge 2\times$ & תיקון מבני\\
\bottomrule
\end{tabular}
\tabcap{כיול מהיר ל־\en{VPN}.}
\end{tablebox}

\begin{tikzpic}
\begin{tikzpicture}[font=\small,node distance=8mm,>={Latex[length=2mm]},
  box/.style={draw,thick,rounded corners,align=center,minimum width=2.4cm,minimum height=1.0cm,fill=cBlueBg}]
  \node[box] (a) {\en{In}};
  \node[box,right=of a] (b) {\en{VPN}};
  \node[box,right=of b,fill=cGreenBg] (c) {\en{Out}};
  \draw[->,thick] (a) -- (b);
  \draw[->,thick] (b) -- (c);
\end{tikzpicture}
\end{tikzpic}
\begin{center}\small איור — זרימה ל־\en{VPN}\end{center}

סדר עבודה:
\begin{enumerate}
\item[(\textbf{א})] מדוד את \en{VPN} לפי~\eqref{eq:tlb}.
\item[(\textbf{ב})] זהה גורם שולט אחד.
\item[(\textbf{ג})] שנה אותו בלבד, ומדוד שוב.
\end{enumerate}

\chapter{זיכרון וירטואלי}
\label{ch:vm}

היחידה נסגרת סביב \en{page}: הגדרה, מדד, ומה נשבר כשמתעלמים ממנו.

\begin{defbox}[frametitle={הגדרה — זיכרון וירטואלי}]
כתובת וירטואלית מפורקת למספר עמוד והיסט בתוך העמוד.
\end{defbox}

\begin{equation}
\label{eq:vm}
VA = VPN \| offset
\end{equation}

המשמעות של~\eqref{eq:vm}: מודדים את \en{page} קודם, ורק אחר כך צוללים לרכיב.

\begin{equation}
\label{eq:vm-check}
\Delta M \propto \Delta x_{\mathrm{dom}}
\end{equation}
כאן $M$ הוא \en{page};~\eqref{eq:vm-check} מזכירה לחפש רגישות, לא מיקרו־אופטימיזציה.

סדר עבודה:
\begin{enumerate}
\item[(\textbf{א})] מדוד את \en{page} לפי~\eqref{eq:vm}.
\item[(\textbf{ב})] זהה גורם שולט אחד.
\item[(\textbf{ג})] שנה אותו בלבד, ומדוד שוב.
\end{enumerate}

\chapter{קוהרנטיות מטמונים}
\label{ch:coherency}

היחידה נסגרת סביב \en{MESI}: הגדרה, מדד, ומה נשבר כשמתעלמים ממנו.

\begin{warnbox}[frametitle={מלכודת — קוהרנטיות מטמונים}]
כתיבה במצב Shared דורשת פסילת עותקים לפני מעבר ל־Modified.
\end{warnbox}

\begin{equation}
\label{eq:coherency}
S \xrightarrow{\mathrm{write}} M
\end{equation}

המשמעות של~\eqref{eq:coherency}: מודדים את \en{MESI} קודם, ורק אחר כך צוללים לרכיב.

סדר עבודה:
\begin{enumerate}
\item[(\textbf{א})] מדוד את \en{MESI} לפי~\eqref{eq:coherency}.
\item[(\textbf{ב})] זהה גורם שולט אחד.
\item[(\textbf{ג})] שנה אותו בלבד, ומדוד שוב.
\end{enumerate}

\chapter{רוחב פס זיכרון}
\label{ch:membw}

היחידה נסגרת סביב \en{bandwidth}: הגדרה, מדד, ומה נשבר כשמתעלמים ממנו.

\begin{keybox}[frametitle={רעיון מפתח — רוחב פס זיכרון}]
רוחב הפס הוא תדר $\times$ רוחב $\times$ ניצולת — לא רק תדר השעון.
\end{keybox}

\begin{equation}
\label{eq:membw}
BW = f\cdot W\cdot U
\end{equation}

המשמעות של~\eqref{eq:membw}: מודדים את \en{bandwidth} קודם, ורק אחר כך צוללים לרכיב.

\begin{equation}
\label{eq:membw-check}
\Delta M \propto \Delta x_{\mathrm{dom}}
\end{equation}
כאן $M$ הוא \en{bandwidth};~\eqref{eq:membw-check} מזכירה לחפש רגישות, לא מיקרו־אופטימיזציה.

רצף שממחיש לחץ על המדד:
\begin{asmblock}
\begin{lstlisting}[style=asm]
lw   $t0, 0($s0)
add  $t1, $t0, $s1   # depends on t0
\end{lstlisting}
\end{asmblock}
התלות על \code{\$t0} היא עלות ב־\en{bandwidth}, לא ״באג קוד״.

סדר עבודה:
\begin{enumerate}
\item[(\textbf{א})] מדוד את \en{bandwidth} לפי~\eqref{eq:membw}.
\item[(\textbf{ב})] זהה גורם שולט אחד.
\item[(\textbf{ג})] שנה אותו בלבד, ומדוד שוב.
\end{enumerate}


\chapter{חיבור — היררכיית זיכרון}
\label{ch:bridge-l1}

הפרקים בחלק ״היררכיית זיכרון״ חולקים משמעת אחת: מדד, נוסחה, מדידה.
הנוסחאות \\eqref{eq:l1}, \\eqref{eq:l2}, \\eqref{eq:repl}, \\eqref{eq:locality} אינן אוסף עובדות — הן חוזה מדידה.

\begin{keybox}[frametitle={חוזה החלק}]
לכל נוסחה ממוספרת בחלק כתוב ניסוי אחד שמאמת אותה. בלי ניסוי — היא דקורציה.
\end{keybox}

כששני מדדים מתנגשים, בוחרים לפי מטרת המערכת (השהיה מול מעבר־נתונים, דיוק מול עלות),
לא לפי הנוסחה היפה יותר.

\part{מערכות הפעלה}

\chapter{תהליכים וחוטים}
\label{ch:proc}

היחידה נסגרת סביב \en{PCB}: הגדרה, מדד, ומה נשבר כשמתעלמים ממנו.

\begin{defbox}[frametitle={הגדרה — תהליכים וחוטים}]
תור המוכנים הוא מה שנשאר אחרי חסימות I/O וסנכרון.
\end{defbox}

\begin{equation}
\label{eq:proc}
N_{\mathrm{ready}} = N - N_{\mathrm{blocked}}
\end{equation}

המשמעות של~\eqref{eq:proc}: מודדים את \en{PCB} קודם, ורק אחר כך צוללים לרכיב.

\begin{tablebox}\centering\small
\begin{tabular}{@{}lll@{}}
\toprule
\textbf{מצב} & \textbf{סדר גודל} & \textbf{פעולה}\\
\midrule
תקין & $\sim 1\times$ & שמור מדידה\\
חשוד & $1.5$–$2\times$ & בודקים תלות\\
שבור & $\ge 2\times$ & תיקון מבני\\
\bottomrule
\end{tabular}
\tabcap{כיול מהיר ל־\en{PCB}.}
\end{tablebox}

סדר עבודה:
\begin{enumerate}
\item[(\textbf{א})] מדוד את \en{PCB} לפי~\eqref{eq:proc}.
\item[(\textbf{ב})] זהה גורם שולט אחד.
\item[(\textbf{ג})] שנה אותו בלבד, ומדוד שוב.
\end{enumerate}

\chapter{תזמון מעבד}
\label{ch:sched}

היחידה נסגרת סביב \en{CFS}: הגדרה, מדד, ומה נשבר כשמתעלמים ממנו.

\begin{keybox}[frametitle={רעיון מפתח — תזמון מעבד}]
\en{CFS} מאזן הוגנות דרך \en{vruntime} משוקלל.
\end{keybox}

\begin{equation}
\label{eq:sched}
vruntime_i += \Delta t \cdot w_0/w_i
\end{equation}

המשמעות של~\eqref{eq:sched}: מודדים את \en{CFS} קודם, ורק אחר כך צוללים לרכיב.

\begin{equation}
\label{eq:sched-check}
\Delta M \propto \Delta x_{\mathrm{dom}}
\end{equation}
כאן $M$ הוא \en{CFS};~\eqref{eq:sched-check} מזכירה לחפש רגישות, לא מיקרו־אופטימיזציה.

סדר עבודה:
\begin{enumerate}
\item[(\textbf{א})] מדוד את \en{CFS} לפי~\eqref{eq:sched}.
\item[(\textbf{ב})] זהה גורם שולט אחד.
\item[(\textbf{ג})] שנה אותו בלבד, ומדוד שוב.
\end{enumerate}

\chapter{סנכרון}
\label{ch:sync}

היחידה נסגרת סביב \en{mutex}: הגדרה, מדד, ומה נשבר כשמתעלמים ממנו.

\begin{warnbox}[frametitle={מלכודת — סנכרון}]
נעילה גלובלית לא משתלבת עם מספר ליבות — צריך פיצול.
\end{warnbox}

\begin{equation}
\label{eq:sync}
P(\mathrm{contention}) \propto N_{\mathrm{cores}}
\end{equation}

המשמעות של~\eqref{eq:sync}: מודדים את \en{mutex} קודם, ורק אחר כך צוללים לרכיב.

\begin{tikzpic}
\begin{tikzpicture}[font=\small,node distance=8mm,>={Latex[length=2mm]},
  box/.style={draw,thick,rounded corners,align=center,minimum width=2.4cm,minimum height=1.0cm,fill=cBlueBg}]
  \node[box] (a) {\en{In}};
  \node[box,right=of a] (b) {\en{mutex}};
  \node[box,right=of b,fill=cGreenBg] (c) {\en{Out}};
  \draw[->,thick] (a) -- (b);
  \draw[->,thick] (b) -- (c);
\end{tikzpicture}
\end{tikzpic}
\begin{center}\small איור — זרימה ל־\en{mutex}\end{center}

סדר עבודה:
\begin{enumerate}
\item[(\textbf{א})] מדוד את \en{mutex} לפי~\eqref{eq:sync}.
\item[(\textbf{ב})] זהה גורם שולט אחד.
\item[(\textbf{ג})] שנה אותו בלבד, ומדוד שוב.
\end{enumerate}

\chapter{קיפאון}
\label{ch:deadlock}

היחידה נסגרת סביב \en{Coffman}: הגדרה, מדד, ומה נשבר כשמתעלמים ממנו.

\begin{thmbox}[frametitle={משפט — קיפאון}]
בלי מעגל בהקצאת משאבים אין קיפאון.
\end{thmbox}

\begin{equation}
\label{eq:deadlock}
\mathrm{safe} \Leftarrow \nexists\,\mathrm{cycle}
\end{equation}

המשמעות של~\eqref{eq:deadlock}: מודדים את \en{Coffman} קודם, ורק אחר כך צוללים לרכיב.

\begin{equation}
\label{eq:deadlock-check}
\Delta M \propto \Delta x_{\mathrm{dom}}
\end{equation}
כאן $M$ הוא \en{Coffman};~\eqref{eq:deadlock-check} מזכירה לחפש רגישות, לא מיקרו־אופטימיזציה.

סדר עבודה:
\begin{enumerate}
\item[(\textbf{א})] מדוד את \en{Coffman} לפי~\eqref{eq:deadlock}.
\item[(\textbf{ב})] זהה גורם שולט אחד.
\item[(\textbf{ג})] שנה אותו בלבד, ומדוד שוב.
\end{enumerate}

\chapter{ניהול זיכרון}
\label{ch:memman}

היחידה נסגרת סביב \en{buddy}: הגדרה, מדד, ומה נשבר כשמתעלמים ממנו.

\begin{defbox}[frametitle={הגדרה — ניהול זיכרון}]
שיטת buddy מקצה בלוקים בחזקות שתיים וממזגת שכנים חופשיים.
\end{defbox}

\begin{equation}
\label{eq:memman}
S_k = 2^k
\end{equation}

המשמעות של~\eqref{eq:memman}: מודדים את \en{buddy} קודם, ורק אחר כך צוללים לרכיב.

\begin{tablebox}\centering\small
\begin{tabular}{@{}lll@{}}
\toprule
\textbf{מצב} & \textbf{סדר גודל} & \textbf{פעולה}\\
\midrule
תקין & $\sim 1\times$ & שמור מדידה\\
חשוד & $1.5$–$2\times$ & בודקים תלות\\
שבור & $\ge 2\times$ & תיקון מבני\\
\bottomrule
\end{tabular}
\tabcap{כיול מהיר ל־\en{buddy}.}
\end{tablebox}

רצף שממחיש לחץ על המדד:
\begin{asmblock}
\begin{lstlisting}[style=asm]
lw   $t0, 0($s0)
add  $t1, $t0, $s1   # depends on t0
\end{lstlisting}
\end{asmblock}
התלות על \code{\$t0} היא עלות ב־\en{buddy}, לא ״באג קוד״.

סדר עבודה:
\begin{enumerate}
\item[(\textbf{א})] מדוד את \en{buddy} לפי~\eqref{eq:memman}.
\item[(\textbf{ב})] זהה גורם שולט אחד.
\item[(\textbf{ג})] שנה אותו בלבד, ומדוד שוב.
\end{enumerate}

\chapter{קלט/פלט}
\label{ch:io}

היחידה נסגרת סביב \en{DMA}: הגדרה, מדד, ומה נשבר כשמתעלמים ממנו.

\begin{keybox}[frametitle={רעיון מפתח — קלט/פלט}]
\en{DMA} מוריד את המעבד ממסלול ההעתקה אחרי \en{setup}.
\end{keybox}

\begin{equation}
\label{eq:io}
T_{\mathrm{IO}}=T_{\mathrm{setup}}+S/BW
\end{equation}

המשמעות של~\eqref{eq:io}: מודדים את \en{DMA} קודם, ורק אחר כך צוללים לרכיב.

\begin{equation}
\label{eq:io-check}
\Delta M \propto \Delta x_{\mathrm{dom}}
\end{equation}
כאן $M$ הוא \en{DMA};~\eqref{eq:io-check} מזכירה לחפש רגישות, לא מיקרו־אופטימיזציה.

סדר עבודה:
\begin{enumerate}
\item[(\textbf{א})] מדוד את \en{DMA} לפי~\eqref{eq:io}.
\item[(\textbf{ב})] זהה גורם שולט אחד.
\item[(\textbf{ג})] שנה אותו בלבד, ומדוד שוב.
\end{enumerate}

\chapter{מערכות קבצים}
\label{ch:fs}

היחידה נסגרת סביב \en{inode}: הגדרה, מדד, ומה נשבר כשמתעלמים ממנו.

\begin{defbox}[frametitle={הגדרה — מערכות קבצים}]
גודל קובץ נגזר ממספר הבלוקים וגודל בלוק.
\end{defbox}

\begin{equation}
\label{eq:fs}
S = n_{\mathrm{blocks}}\cdot B
\end{equation}

המשמעות של~\eqref{eq:fs}: מודדים את \en{inode} קודם, ורק אחר כך צוללים לרכיב.

סדר עבודה:
\begin{enumerate}
\item[(\textbf{א})] מדוד את \en{inode} לפי~\eqref{eq:fs}.
\item[(\textbf{ב})] זהה גורם שולט אחד.
\item[(\textbf{ג})] שנה אותו בלבד, ומדוד שוב.
\end{enumerate}

\chapter{וירטואליזציה}
\label{ch:virt}

היחידה נסגרת סביב \en{VMEXIT}: הגדרה, מדד, ומה נשבר כשמתעלמים ממנו.

\begin{warnbox}[frametitle={מלכודת — וירטואליזציה}]
תדירות \en{VMEXIT} גבוהה הורסת את יתרון הווירטואליזציה.
\end{warnbox}

\begin{equation}
\label{eq:virt}
T_{\mathrm{virt}}=T_{\mathrm{guest}}+f_{\mathrm{exit}}T_{\mathrm{exit}}
\end{equation}

המשמעות של~\eqref{eq:virt}: מודדים את \en{VMEXIT} קודם, ורק אחר כך צוללים לרכיב.

\begin{equation}
\label{eq:virt-check}
\Delta M \propto \Delta x_{\mathrm{dom}}
\end{equation}
כאן $M$ הוא \en{VMEXIT};~\eqref{eq:virt-check} מזכירה לחפש רגישות, לא מיקרו־אופטימיזציה.

סדר עבודה:
\begin{enumerate}
\item[(\textbf{א})] מדוד את \en{VMEXIT} לפי~\eqref{eq:virt}.
\item[(\textbf{ב})] זהה גורם שולט אחד.
\item[(\textbf{ג})] שנה אותו בלבד, ומדוד שוב.
\end{enumerate}


\chapter{חיבור — מערכות הפעלה}
\label{ch:bridge-proc}

הפרקים בחלק ״מערכות הפעלה״ חולקים משמעת אחת: מדד, נוסחה, מדידה.
הנוסחאות \\eqref{eq:proc}, \\eqref{eq:sched}, \\eqref{eq:sync}, \\eqref{eq:deadlock} אינן אוסף עובדות — הן חוזה מדידה.

\begin{keybox}[frametitle={חוזה החלק}]
לכל נוסחה ממוספרת בחלק כתוב ניסוי אחד שמאמת אותה. בלי ניסוי — היא דקורציה.
\end{keybox}

כששני מדדים מתנגשים, בוחרים לפי מטרת המערכת (השהיה מול מעבר־נתונים, דיוק מול עלות),
לא לפי הנוסחה היפה יותר.

\part{רשתות ותקשורת}

\chapter{מודל השכבות}
\label{ch:layers}

היחידה נסגרת סביב \en{OSI}: הגדרה, מדד, ומה נשבר כשמתעלמים ממנו.

\begin{defbox}[frametitle={הגדרה — מודל השכבות}]
כל שכבה מוסיפה כותרת ומעבירה מטה — חוזה ברור בין שכבות.
\end{defbox}

\begin{equation}
\label{eq:layers}
L_i \to L_{i-1}
\end{equation}

המשמעות של~\eqref{eq:layers}: מודדים את \en{OSI} קודם, ורק אחר כך צוללים לרכיב.

\begin{tablebox}\centering\small
\begin{tabular}{@{}lll@{}}
\toprule
\textbf{מצב} & \textbf{סדר גודל} & \textbf{פעולה}\\
\midrule
תקין & $\sim 1\times$ & שמור מדידה\\
חשוד & $1.5$–$2\times$ & בודקים תלות\\
שבור & $\ge 2\times$ & תיקון מבני\\
\bottomrule
\end{tabular}
\tabcap{כיול מהיר ל־\en{OSI}.}
\end{tablebox}

\begin{tikzpic}
\begin{tikzpicture}[font=\small,node distance=8mm,>={Latex[length=2mm]},
  box/.style={draw,thick,rounded corners,align=center,minimum width=2.4cm,minimum height=1.0cm,fill=cBlueBg}]
  \node[box] (a) {\en{In}};
  \node[box,right=of a] (b) {\en{OSI}};
  \node[box,right=of b,fill=cGreenBg] (c) {\en{Out}};
  \draw[->,thick] (a) -- (b);
  \draw[->,thick] (b) -- (c);
\end{tikzpicture}
\end{tikzpic}
\begin{center}\small איור — זרימה ל־\en{OSI}\end{center}

סדר עבודה:
\begin{enumerate}
\item[(\textbf{א})] מדוד את \en{OSI} לפי~\eqref{eq:layers}.
\item[(\textbf{ב})] זהה גורם שולט אחד.
\item[(\textbf{ג})] שנה אותו בלבד, ומדוד שוב.
\end{enumerate}

\chapter{TCP}
\label{ch:tcp}

היחידה נסגרת סביב \en{cwnd}: הגדרה, מדד, ומה נשבר כשמתעלמים ממנו.

\begin{keybox}[frametitle={רעיון מפתח — TCP}]
העלאת \en{cwnd} ב־\en{AIMD} היא זהירה אחרי \en{slow start}.
\end{keybox}

\begin{equation}
\label{eq:tcp}
cwnd_{t+1}=cwnd_t + 1/\mathrm{cwnd}_t
\end{equation}

המשמעות של~\eqref{eq:tcp}: מודדים את \en{cwnd} קודם, ורק אחר כך צוללים לרכיב.

\begin{equation}
\label{eq:tcp-check}
\Delta M \propto \Delta x_{\mathrm{dom}}
\end{equation}
כאן $M$ הוא \en{cwnd};~\eqref{eq:tcp-check} מזכירה לחפש רגישות, לא מיקרו־אופטימיזציה.

רצף שממחיש לחץ על המדד:
\begin{asmblock}
\begin{lstlisting}[style=asm]
lw   $t0, 0($s0)
add  $t1, $t0, $s1   # depends on t0
\end{lstlisting}
\end{asmblock}
התלות על \code{\$t0} היא עלות ב־\en{cwnd}, לא ״באג קוד״.

סדר עבודה:
\begin{enumerate}
\item[(\textbf{א})] מדוד את \en{cwnd} לפי~\eqref{eq:tcp}.
\item[(\textbf{ב})] זהה גורם שולט אחד.
\item[(\textbf{ג})] שנה אותו בלבד, ומדוד שוב.
\end{enumerate}

\chapter{ניתוב}
\label{ch:routing}

היחידה נסגרת סביב \en{SPF}: הגדרה, מדד, ומה נשבר כשמתעלמים ממנו.

\begin{defbox}[frametitle={הגדרה — ניתוב}]
דייקסטרה בונה מרחקים קצרים ביותר מעץ שורש.
\end{defbox}

\begin{equation}
\label{eq:routing}
d(v)=\min_u\, d(u)+w(u,v)
\end{equation}

המשמעות של~\eqref{eq:routing}: מודדים את \en{SPF} קודם, ורק אחר כך צוללים לרכיב.

סדר עבודה:
\begin{enumerate}
\item[(\textbf{א})] מדוד את \en{SPF} לפי~\eqref{eq:routing}.
\item[(\textbf{ב})] זהה גורם שולט אחד.
\item[(\textbf{ג})] שנה אותו בלבד, ומדוד שוב.
\end{enumerate}

\chapter{עומסים ו־QoS}
\label{ch:qos}

היחידה נסגרת סביב \en{latency}: הגדרה, מדד, ומה נשבר כשמתעלמים ממנו.

\begin{keybox}[frametitle={רעיון מפתח — עומסים ו־QoS}]
השהיה היא סכום התפשטות, תור ושידור — לא רק רוחב פס.
\end{keybox}

\begin{equation}
\label{eq:qos}
L = L_{\mathrm{prop}}+L_{\mathrm{queue}}+L_{\mathrm{tx}}
\end{equation}

המשמעות של~\eqref{eq:qos}: מודדים את \en{latency} קודם, ורק אחר כך צוללים לרכיב.

\begin{equation}
\label{eq:qos-check}
\Delta M \propto \Delta x_{\mathrm{dom}}
\end{equation}
כאן $M$ הוא \en{latency};~\eqref{eq:qos-check} מזכירה לחפש רגישות, לא מיקרו־אופטימיזציה.

סדר עבודה:
\begin{enumerate}
\item[(\textbf{א})] מדוד את \en{latency} לפי~\eqref{eq:qos}.
\item[(\textbf{ב})] זהה גורם שולט אחד.
\item[(\textbf{ג})] שנה אותו בלבד, ומדוד שוב.
\end{enumerate}

\chapter{אבטחת רשת}
\label{ch:netsec}

היחידה נסגרת סביב \en{TLS}: הגדרה, מדד, ומה נשבר כשמתעלמים ממנו.

\begin{defbox}[frametitle={הגדרה — אבטחת רשת}]
\en{TLS} מפריד משא ומתן מפתחות מרישום הנתונים המוצפן.
\end{defbox}

\begin{equation}
\label{eq:netsec}
\mathrm{TLS}=\mathrm{handshake}+\mathrm{record}
\end{equation}

המשמעות של~\eqref{eq:netsec}: מודדים את \en{TLS} קודם, ורק אחר כך צוללים לרכיב.

\begin{tablebox}\centering\small
\begin{tabular}{@{}lll@{}}
\toprule
\textbf{מצב} & \textbf{סדר גודל} & \textbf{פעולה}\\
\midrule
תקין & $\sim 1\times$ & שמור מדידה\\
חשוד & $1.5$–$2\times$ & בודקים תלות\\
שבור & $\ge 2\times$ & תיקון מבני\\
\bottomrule
\end{tabular}
\tabcap{כיול מהיר ל־\en{TLS}.}
\end{tablebox}

סדר עבודה:
\begin{enumerate}
\item[(\textbf{א})] מדוד את \en{TLS} לפי~\eqref{eq:netsec}.
\item[(\textbf{ב})] זהה גורם שולט אחד.
\item[(\textbf{ג})] שנה אותו בלבד, ומדוד שוב.
\end{enumerate}

\chapter{CDN ומטמונים}
\label{ch:cdn}

היחידה נסגרת סביב \en{HIT}: הגדרה, מדד, ומה נשבר כשמתעלמים ממנו.

\begin{keybox}[frametitle={רעיון מפתח — CDN ומטמונים}]
\en{CDN} הוא \en{AMAT} לרשת: פגיעה בקצה מול החמצה למקור.
\end{keybox}

\begin{equation}
\label{eq:cdn}
T_{\mathrm{user}}=T_{\mathrm{edge}}+r_{\mathrm{miss}}T_{\mathrm{origin}}
\end{equation}

המשמעות של~\eqref{eq:cdn}: מודדים את \en{HIT} קודם, ורק אחר כך צוללים לרכיב.

\begin{equation}
\label{eq:cdn-check}
\Delta M \propto \Delta x_{\mathrm{dom}}
\end{equation}
כאן $M$ הוא \en{HIT};~\eqref{eq:cdn-check} מזכירה לחפש רגישות, לא מיקרו־אופטימיזציה.

סדר עבודה:
\begin{enumerate}
\item[(\textbf{א})] מדוד את \en{HIT} לפי~\eqref{eq:cdn}.
\item[(\textbf{ב})] זהה גורם שולט אחד.
\item[(\textbf{ג})] שנה אותו בלבד, ומדוד שוב.
\end{enumerate}


\chapter{חיבור — רשתות ותקשורת}
\label{ch:bridge-layers}

הפרקים בחלק ״רשתות ותקשורת״ חולקים משמעת אחת: מדד, נוסחה, מדידה.
הנוסחאות \\eqref{eq:layers}, \\eqref{eq:tcp}, \\eqref{eq:routing}, \\eqref{eq:qos} אינן אוסף עובדות — הן חוזה מדידה.

\begin{keybox}[frametitle={חוזה החלק}]
לכל נוסחה ממוספרת בחלק כתוב ניסוי אחד שמאמת אותה. בלי ניסוי — היא דקורציה.
\end{keybox}

כששני מדדים מתנגשים, בוחרים לפי מטרת המערכת (השהיה מול מעבר־נתונים, דיוק מול עלות),
לא לפי הנוסחה היפה יותר.

\part{אלגוריתמים ומבני נתונים}

\chapter{סיבוכיות}
\label{ch:complexity}

היחידה נסגרת סביב \en{BigO}: הגדרה, מדד, ומה נשבר כשמתעלמים ממנו.

\begin{defbox}[frametitle={הגדרה — סיבוכיות}]
סימון אסימפטוטי מתעלם מקבועים ומתמקד בסדר הגדילה.
\end{defbox}

\begin{equation}
\label{eq:complexity}
T(n)=\Theta(n\log n)
\end{equation}

המשמעות של~\eqref{eq:complexity}: מודדים את \en{BigO} קודם, ורק אחר כך צוללים לרכיב.

רצף שממחיש לחץ על המדד:
\begin{asmblock}
\begin{lstlisting}[style=asm]
lw   $t0, 0($s0)
add  $t1, $t0, $s1   # depends on t0
\end{lstlisting}
\end{asmblock}
התלות על \code{\$t0} היא עלות ב־\en{BigO}, לא ״באג קוד״.

\begin{tikzpic}
\begin{tikzpicture}[font=\small,node distance=8mm,>={Latex[length=2mm]},
  box/.style={draw,thick,rounded corners,align=center,minimum width=2.4cm,minimum height=1.0cm,fill=cBlueBg}]
  \node[box] (a) {\en{In}};
  \node[box,right=of a] (b) {\en{BigO}};
  \node[box,right=of b,fill=cGreenBg] (c) {\en{Out}};
  \draw[->,thick] (a) -- (b);
  \draw[->,thick] (b) -- (c);
\end{tikzpicture}
\end{tikzpic}
\begin{center}\small איור — זרימה ל־\en{BigO}\end{center}

סדר עבודה:
\begin{enumerate}
\item[(\textbf{א})] מדוד את \en{BigO} לפי~\eqref{eq:complexity}.
\item[(\textbf{ב})] זהה גורם שולט אחד.
\item[(\textbf{ג})] שנה אותו בלבד, ומדוד שוב.
\end{enumerate}

\chapter{מיונים}
\label{ch:sort}

היחידה נסגרת סביב \en{quicksort}: הגדרה, מדד, ומה נשבר כשמתעלמים ממנו.

\begin{keybox}[frametitle={רעיון מפתח — מיונים}]
פיצול מאוזן נותן $n\log n$; פיצול גרוע מחזיר $n^2$.
\end{keybox}

\begin{equation}
\label{eq:sort}
T(n)=2T(n/2)+\Theta(n)
\end{equation}

המשמעות של~\eqref{eq:sort}: מודדים את \en{quicksort} קודם, ורק אחר כך צוללים לרכיב.

\begin{equation}
\label{eq:sort-check}
\Delta M \propto \Delta x_{\mathrm{dom}}
\end{equation}
כאן $M$ הוא \en{quicksort};~\eqref{eq:sort-check} מזכירה לחפש רגישות, לא מיקרו־אופטימיזציה.

סדר עבודה:
\begin{enumerate}
\item[(\textbf{א})] מדוד את \en{quicksort} לפי~\eqref{eq:sort}.
\item[(\textbf{ב})] זהה גורם שולט אחד.
\item[(\textbf{ג})] שנה אותו בלבד, ומדוד שוב.
\end{enumerate}

\chapter{גרפים}
\label{ch:graphs}

היחידה נסגרת סביב \en{BFS}: הגדרה, מדד, ומה נשבר כשמתעלמים ממנו.

\begin{defbox}[frametitle={הגדרה — גרפים}]
עץ פורש על $|V|$ קודקודים מכיל בדיוק $|V|-1$ קשתות.
\end{defbox}

\begin{equation}
\label{eq:graphs}
|E|_{\mathrm{tree}}=|V|-1
\end{equation}

המשמעות של~\eqref{eq:graphs}: מודדים את \en{BFS} קודם, ורק אחר כך צוללים לרכיב.

\begin{tablebox}\centering\small
\begin{tabular}{@{}lll@{}}
\toprule
\textbf{מצב} & \textbf{סדר גודל} & \textbf{פעולה}\\
\midrule
תקין & $\sim 1\times$ & שמור מדידה\\
חשוד & $1.5$–$2\times$ & בודקים תלות\\
שבור & $\ge 2\times$ & תיקון מבני\\
\bottomrule
\end{tabular}
\tabcap{כיול מהיר ל־\en{BFS}.}
\end{tablebox}

סדר עבודה:
\begin{enumerate}
\item[(\textbf{א})] מדוד את \en{BFS} לפי~\eqref{eq:graphs}.
\item[(\textbf{ב})] זהה גורם שולט אחד.
\item[(\textbf{ג})] שנה אותו בלבד, ומדוד שוב.
\end{enumerate}

\chapter{תכנון דינמי}
\label{ch:dp}

היחידה נסגרת סביב \en{opt}: הגדרה, מדד, ומה נשבר כשמתעלמים ממנו.

\begin{keybox}[frametitle={רעיון מפתח — תכנון דינמי}]
אופטימום נבנה מאופטימומים של תתי־בעיות חופפות.
\end{keybox}

\begin{equation}
\label{eq:dp}
OPT(i)=\min_j OPT(j)+c(j,i)
\end{equation}

המשמעות של~\eqref{eq:dp}: מודדים את \en{opt} קודם, ורק אחר כך צוללים לרכיב.

\begin{equation}
\label{eq:dp-check}
\Delta M \propto \Delta x_{\mathrm{dom}}
\end{equation}
כאן $M$ הוא \en{opt};~\eqref{eq:dp-check} מזכירה לחפש רגישות, לא מיקרו־אופטימיזציה.

סדר עבודה:
\begin{enumerate}
\item[(\textbf{א})] מדוד את \en{opt} לפי~\eqref{eq:dp}.
\item[(\textbf{ב})] זהה גורם שולט אחד.
\item[(\textbf{ג})] שנה אותו בלבד, ומדוד שוב.
\end{enumerate}

\chapter{גיבוב}
\label{ch:hash}

היחידה נסגרת סביב \en{chaining}: הגדרה, מדד, ומה נשבר כשמתעלמים ממנו.

\begin{defbox}[frametitle={הגדרה — גיבוב}]
מקדם העומס $\alpha$ שולט באורך השרשראות הממוצע.
\end{defbox}

\begin{equation}
\label{eq:hash}
\alpha = n/m
\end{equation}

המשמעות של~\eqref{eq:hash}: מודדים את \en{chaining} קודם, ורק אחר כך צוללים לרכיב.

סדר עבודה:
\begin{enumerate}
\item[(\textbf{א})] מדוד את \en{chaining} לפי~\eqref{eq:hash}.
\item[(\textbf{ב})] זהה גורם שולט אחד.
\item[(\textbf{ג})] שנה אותו בלבד, ומדוד שוב.
\end{enumerate}

\chapter{עצים מאוזנים}
\label{ch:trees}

היחידה נסגרת סביב \en{AVL}: הגדרה, מדד, ומה נשבר כשמתעלמים ממנו.

\begin{thmbox}[frametitle={משפט — עצים מאוזנים}]
איזון גבהים מבטיח חיפוש לוגריתמי.
\end{thmbox}

\begin{equation}
\label{eq:trees}
h = O(\log n)
\end{equation}

המשמעות של~\eqref{eq:trees}: מודדים את \en{AVL} קודם, ורק אחר כך צוללים לרכיב.

\begin{equation}
\label{eq:trees-check}
\Delta M \propto \Delta x_{\mathrm{dom}}
\end{equation}
כאן $M$ הוא \en{AVL};~\eqref{eq:trees-check} מזכירה לחפש רגישות, לא מיקרו־אופטימיזציה.

רצף שממחיש לחץ על המדד:
\begin{asmblock}
\begin{lstlisting}[style=asm]
lw   $t0, 0($s0)
add  $t1, $t0, $s1   # depends on t0
\end{lstlisting}
\end{asmblock}
התלות על \code{\$t0} היא עלות ב־\en{AVL}, לא ״באג קוד״.

סדר עבודה:
\begin{enumerate}
\item[(\textbf{א})] מדוד את \en{AVL} לפי~\eqref{eq:trees}.
\item[(\textbf{ב})] זהה גורם שולט אחד.
\item[(\textbf{ג})] שנה אותו בלבד, ומדוד שוב.
\end{enumerate}

\chapter{זרימה ברשתות}
\label{ch:flow}

היחידה נסגרת סביב \en{maxflow}: הגדרה, מדד, ומה נשבר כשמתעלמים ממנו.

\begin{defbox}[frametitle={הגדרה — זרימה ברשתות}]
ערך הזרימה הוא סכום היציאות מהמקור.
\end{defbox}

\begin{equation}
\label{eq:flow}
v(f)=\sum_{e \ni s} f(e)
\end{equation}

המשמעות של~\eqref{eq:flow}: מודדים את \en{maxflow} קודם, ורק אחר כך צוללים לרכיב.

\begin{tablebox}\centering\small
\begin{tabular}{@{}lll@{}}
\toprule
\textbf{מצב} & \textbf{סדר גודל} & \textbf{פעולה}\\
\midrule
תקין & $\sim 1\times$ & שמור מדידה\\
חשוד & $1.5$–$2\times$ & בודקים תלות\\
שבור & $\ge 2\times$ & תיקון מבני\\
\bottomrule
\end{tabular}
\tabcap{כיול מהיר ל־\en{maxflow}.}
\end{tablebox}

\begin{tikzpic}
\begin{tikzpicture}[font=\small,node distance=8mm,>={Latex[length=2mm]},
  box/.style={draw,thick,rounded corners,align=center,minimum width=2.4cm,minimum height=1.0cm,fill=cBlueBg}]
  \node[box] (a) {\en{In}};
  \node[box,right=of a] (b) {\en{maxflow}};
  \node[box,right=of b,fill=cGreenBg] (c) {\en{Out}};
  \draw[->,thick] (a) -- (b);
  \draw[->,thick] (b) -- (c);
\end{tikzpicture}
\end{tikzpic}
\begin{center}\small איור — זרימה ל־\en{maxflow}\end{center}

סדר עבודה:
\begin{enumerate}
\item[(\textbf{א})] מדוד את \en{maxflow} לפי~\eqref{eq:flow}.
\item[(\textbf{ב})] זהה גורם שולט אחד.
\item[(\textbf{ג})] שנה אותו בלבד, ומדוד שוב.
\end{enumerate}

\chapter{קירובים}
\label{ch:approx}

היחידה נסגרת סביב \en{ratio}: הגדרה, מדד, ומה נשבר כשמתעלמים ממנו.

\begin{keybox}[frametitle={רעיון מפתח — קירובים}]
יחס הקירוב $\rho$ הוא ההבטחה כש־\en{NP}-קשה מונע דיוק.
\end{keybox}

\begin{equation}
\label{eq:approx}
ALG \le \rho\cdot OPT
\end{equation}

המשמעות של~\eqref{eq:approx}: מודדים את \en{ratio} קודם, ורק אחר כך צוללים לרכיב.

\begin{equation}
\label{eq:approx-check}
\Delta M \propto \Delta x_{\mathrm{dom}}
\end{equation}
כאן $M$ הוא \en{ratio};~\eqref{eq:approx-check} מזכירה לחפש רגישות, לא מיקרו־אופטימיזציה.

סדר עבודה:
\begin{enumerate}
\item[(\textbf{א})] מדוד את \en{ratio} לפי~\eqref{eq:approx}.
\item[(\textbf{ב})] זהה גורם שולט אחד.
\item[(\textbf{ג})] שנה אותו בלבד, ומדוד שוב.
\end{enumerate}


\chapter{חיבור — אלגוריתמים ומבני נתונים}
\label{ch:bridge-complexity}

הפרקים בחלק ״אלגוריתמים ומבני נתונים״ חולקים משמעת אחת: מדד, נוסחה, מדידה.
הנוסחאות \\eqref{eq:complexity}, \\eqref{eq:sort}, \\eqref{eq:graphs}, \\eqref{eq:dp} אינן אוסף עובדות — הן חוזה מדידה.

\begin{keybox}[frametitle={חוזה החלק}]
לכל נוסחה ממוספרת בחלק כתוב ניסוי אחד שמאמת אותה. בלי ניסוי — היא דקורציה.
\end{keybox}

כששני מדדים מתנגשים, בוחרים לפי מטרת המערכת (השהיה מול מעבר־נתונים, דיוק מול עלות),
לא לפי הנוסחה היפה יותר.

\part{למידת מכונה}

\chapter{רגרסיה לינארית}
\label{ch:linreg}

היחידה נסגרת סביב \en{MSE}: הגדרה, מדד, ומה נשבר כשמתעלמים ממנו.

\begin{defbox}[frametitle={הגדרה — רגרסיה לינארית}]
פחות ריבועים נותן את האומדן הליניארי בצורת מטריצה סגורה.
\end{defbox}

\begin{equation}
\label{eq:linreg}
\hat\beta=(X^\top X)^{-1}X^\top y
\end{equation}

המשמעות של~\eqref{eq:linreg}: מודדים את \en{MSE} קודם, ורק אחר כך צוללים לרכיב.

סדר עבודה:
\begin{enumerate}
\item[(\textbf{א})] מדוד את \en{MSE} לפי~\eqref{eq:linreg}.
\item[(\textbf{ב})] זהה גורם שולט אחד.
\item[(\textbf{ג})] שנה אותו בלבד, ומדוד שוב.
\end{enumerate}

\chapter{סיווג}
\label{ch:clf}

היחידה נסגרת סביב \en{logistic}: הגדרה, מדד, ומה נשבר כשמתעלמים ממנו.

\begin{keybox}[frametitle={רעיון מפתח — סיווג}]
לוגיסטי ממפה ציון ליניארי להסתברות דרך $\sigma$.
\end{keybox}

\begin{equation}
\label{eq:clf}
p=\sigma(w^\top x)
\end{equation}

המשמעות של~\eqref{eq:clf}: מודדים את \en{logistic} קודם, ורק אחר כך צוללים לרכיב.

\begin{equation}
\label{eq:clf-check}
\Delta M \propto \Delta x_{\mathrm{dom}}
\end{equation}
כאן $M$ הוא \en{logistic};~\eqref{eq:clf-check} מזכירה לחפש רגישות, לא מיקרו־אופטימיזציה.

סדר עבודה:
\begin{enumerate}
\item[(\textbf{א})] מדוד את \en{logistic} לפי~\eqref{eq:clf}.
\item[(\textbf{ב})] זהה גורם שולט אחד.
\item[(\textbf{ג})] שנה אותו בלבד, ומדוד שוב.
\end{enumerate}

\chapter{רגולריזציה}
\label{ch:reg}

היחידה נסגרת סביב \en{ridge}: הגדרה, מדד, ומה נשבר כשמתעלמים ממנו.

\begin{warnbox}[frametitle={מלכודת — רגולריזציה}]
בלי $\lambda$ המודל משנן רעש; עם $\lambda$ גדול מדי הוא מפספס אות.
\end{warnbox}

\begin{equation}
\label{eq:reg}
\mathcal{L}=\|y-Xw\|^2+\lambda\|w\|^2
\end{equation}

המשמעות של~\eqref{eq:reg}: מודדים את \en{ridge} קודם, ורק אחר כך צוללים לרכיב.

\begin{tablebox}\centering\small
\begin{tabular}{@{}lll@{}}
\toprule
\textbf{מצב} & \textbf{סדר גודל} & \textbf{פעולה}\\
\midrule
תקין & $\sim 1\times$ & שמור מדידה\\
חשוד & $1.5$–$2\times$ & בודקים תלות\\
שבור & $\ge 2\times$ & תיקון מבני\\
\bottomrule
\end{tabular}
\tabcap{כיול מהיר ל־\en{ridge}.}
\end{tablebox}

רצף שממחיש לחץ על המדד:
\begin{asmblock}
\begin{lstlisting}[style=asm]
lw   $t0, 0($s0)
add  $t1, $t0, $s1   # depends on t0
\end{lstlisting}
\end{asmblock}
התלות על \code{\$t0} היא עלות ב־\en{ridge}, לא ״באג קוד״.

סדר עבודה:
\begin{enumerate}
\item[(\textbf{א})] מדוד את \en{ridge} לפי~\eqref{eq:reg}.
\item[(\textbf{ב})] זהה גורם שולט אחד.
\item[(\textbf{ג})] שנה אותו בלבד, ומדוד שוב.
\end{enumerate}

\chapter{רשתות עצביות}
\label{ch:nn}

היחידה נסגרת סביב \en{backprop}: הגדרה, מדד, ומה נשבר כשמתעלמים ממנו.

\begin{defbox}[frametitle={הגדרה — רשתות עצביות}]
שגיאה מתפשטת אחורה דרך היעקוביאן של כל שכבה.
\end{defbox}

\begin{equation}
\label{eq:nn}
\delta^{(l)}=(W^{(l+1)})^\top\delta^{(l+1)}\odot\sigma'(z^{(l)})
\end{equation}

המשמעות של~\eqref{eq:nn}: מודדים את \en{backprop} קודם, ורק אחר כך צוללים לרכיב.

\begin{equation}
\label{eq:nn-check}
\Delta M \propto \Delta x_{\mathrm{dom}}
\end{equation}
כאן $M$ הוא \en{backprop};~\eqref{eq:nn-check} מזכירה לחפש רגישות, לא מיקרו־אופטימיזציה.

סדר עבודה:
\begin{enumerate}
\item[(\textbf{א})] מדוד את \en{backprop} לפי~\eqref{eq:nn}.
\item[(\textbf{ב})] זהה גורם שולט אחד.
\item[(\textbf{ג})] שנה אותו בלבד, ומדוד שוב.
\end{enumerate}

\chapter{הכללה}
\label{ch:gen}

היחידה נסגרת סביב \en{VC}: הגדרה, מדד, ומה נשבר כשמתעלמים ממנו.

\begin{thmbox}[frametitle={משפט — הכללה}]
פער הכללה קטן כש־$n$ גדול ביחס לסיבוכיות $d$.
\end{thmbox}

\begin{equation}
\label{eq:gen}
R(h)\le\hat R(h)+O\!\left(\sqrt{d/n}\right)
\end{equation}

המשמעות של~\eqref{eq:gen}: מודדים את \en{VC} קודם, ורק אחר כך צוללים לרכיב.

\begin{tikzpic}
\begin{tikzpicture}[font=\small,node distance=8mm,>={Latex[length=2mm]},
  box/.style={draw,thick,rounded corners,align=center,minimum width=2.4cm,minimum height=1.0cm,fill=cBlueBg}]
  \node[box] (a) {\en{In}};
  \node[box,right=of a] (b) {\en{VC}};
  \node[box,right=of b,fill=cGreenBg] (c) {\en{Out}};
  \draw[->,thick] (a) -- (b);
  \draw[->,thick] (b) -- (c);
\end{tikzpicture}
\end{tikzpic}
\begin{center}\small איור — זרימה ל־\en{VC}\end{center}

סדר עבודה:
\begin{enumerate}
\item[(\textbf{א})] מדוד את \en{VC} לפי~\eqref{eq:gen}.
\item[(\textbf{ב})] זהה גורם שולט אחד.
\item[(\textbf{ג})] שנה אותו בלבד, ומדוד שוב.
\end{enumerate}

\chapter{למידה לא מפוקחת}
\label{ch:unsup}

היחידה נסגרת סביב \en{kmeans}: הגדרה, מדד, ומה נשבר כשמתעלמים ממנו.

\begin{keybox}[frametitle={רעיון מפתח — למידה לא מפוקחת}]
\en{k-means} ממזער ריבועי מרחק למרכזים — תלוי באתחול.
\end{keybox}

\begin{equation}
\label{eq:unsup}
J=\sum_k\sum_{x\in C_k}\|x-\mu_k\|^2
\end{equation}

המשמעות של~\eqref{eq:unsup}: מודדים את \en{kmeans} קודם, ורק אחר כך צוללים לרכיב.

\begin{equation}
\label{eq:unsup-check}
\Delta M \propto \Delta x_{\mathrm{dom}}
\end{equation}
כאן $M$ הוא \en{kmeans};~\eqref{eq:unsup-check} מזכירה לחפש רגישות, לא מיקרו־אופטימיזציה.

סדר עבודה:
\begin{enumerate}
\item[(\textbf{א})] מדוד את \en{kmeans} לפי~\eqref{eq:unsup}.
\item[(\textbf{ב})] זהה גורם שולט אחד.
\item[(\textbf{ג})] שנה אותו בלבד, ומדוד שוב.
\end{enumerate}

\chapter{הטמעות}
\label{ch:embed}

היחידה נסגרת סביב \en{word2vec}: הגדרה, מדד, ומה נשבר כשמתעלמים ממנו.

\begin{defbox}[frametitle={הגדרה — הטמעות}]
הטמעה לומדת הקשר מקומי: מילים קרובות מקבלות וקטורים קרובים.
\end{defbox}

\begin{equation}
\label{eq:embed}
\max \sum_t \log p(w_{t+j}\mid w_t)
\end{equation}

המשמעות של~\eqref{eq:embed}: מודדים את \en{word2vec} קודם, ורק אחר כך צוללים לרכיב.

\begin{tablebox}\centering\small
\begin{tabular}{@{}lll@{}}
\toprule
\textbf{מצב} & \textbf{סדר גודל} & \textbf{פעולה}\\
\midrule
תקין & $\sim 1\times$ & שמור מדידה\\
חשוד & $1.5$–$2\times$ & בודקים תלות\\
שבור & $\ge 2\times$ & תיקון מבני\\
\bottomrule
\end{tabular}
\tabcap{כיול מהיר ל־\en{word2vec}.}
\end{tablebox}

סדר עבודה:
\begin{enumerate}
\item[(\textbf{א})] מדוד את \en{word2vec} לפי~\eqref{eq:embed}.
\item[(\textbf{ב})] זהה גורם שולט אחד.
\item[(\textbf{ג})] שנה אותו בלבד, ומדוד שוב.
\end{enumerate}


\chapter{חיבור — למידת מכונה}
\label{ch:bridge-linreg}

הפרקים בחלק ״למידת מכונה״ חולקים משמעת אחת: מדד, נוסחה, מדידה.
הנוסחאות \\eqref{eq:linreg}, \\eqref{eq:clf}, \\eqref{eq:reg}, \\eqref{eq:nn} אינן אוסף עובדות — הן חוזה מדידה.

\begin{keybox}[frametitle={חוזה החלק}]
לכל נוסחה ממוספרת בחלק כתוב ניסוי אחד שמאמת אותה. בלי ניסוי — היא דקורציה.
\end{keybox}

כששני מדדים מתנגשים, בוחרים לפי מטרת המערכת (השהיה מול מעבר־נתונים, דיוק מול עלות),
לא לפי הנוסחה היפה יותר.

\part{חישוב קוונטי}

\chapter{הקיוביט}
\label{ch:qubit}

היחידה נסגרת סביב \en{Bloch}: הגדרה, מדד, ומה נשבר כשמתעלמים ממנו.

\begin{defbox}[frametitle={הגדרה — הקיוביט}]
מצב טהור הוא נקודה על כדור בלוך — שני פרמטרים ממשיים.
\end{defbox}

\begin{equation}
\label{eq:qubit}
\ket\psi=\cos\tfrac\theta2\ket0+e^{i\varphi}\sin\tfrac\theta2\ket1
\end{equation}

המשמעות של~\eqref{eq:qubit}: מודדים את \en{Bloch} קודם, ורק אחר כך צוללים לרכיב.

\begin{equation}
\label{eq:qubit-check}
\Delta M \propto \Delta x_{\mathrm{dom}}
\end{equation}
כאן $M$ הוא \en{Bloch};~\eqref{eq:qubit-check} מזכירה לחפש רגישות, לא מיקרו־אופטימיזציה.

רצף שממחיש לחץ על המדד:
\begin{asmblock}
\begin{lstlisting}[style=asm]
lw   $t0, 0($s0)
add  $t1, $t0, $s1   # depends on t0
\end{lstlisting}
\end{asmblock}
התלות על \code{\$t0} היא עלות ב־\en{Bloch}, לא ״באג קוד״.

סדר עבודה:
\begin{enumerate}
\item[(\textbf{א})] מדוד את \en{Bloch} לפי~\eqref{eq:qubit}.
\item[(\textbf{ב})] זהה גורם שולט אחד.
\item[(\textbf{ג})] שנה אותו בלבד, ומדוד שוב.
\end{enumerate}

\chapter{שערים בסיסיים}
\label{ch:gates}

היחידה נסגרת סביב \en{Hadamard}: הגדרה, מדד, ומה נשבר כשמתעלמים ממנו.

\begin{keybox}[frametitle={רעיון מפתח — שערים בסיסיים}]
הדמרד יוצר סופרפוזיציה שווה מ־$\ket0$.
\end{keybox}

\begin{equation}
\label{eq:gates}
H\ket0=\ket{+}
\end{equation}

המשמעות של~\eqref{eq:gates}: מודדים את \en{Hadamard} קודם, ורק אחר כך צוללים לרכיב.

סדר עבודה:
\begin{enumerate}
\item[(\textbf{א})] מדוד את \en{Hadamard} לפי~\eqref{eq:gates}.
\item[(\textbf{ב})] זהה גורם שולט אחד.
\item[(\textbf{ג})] שנה אותו בלבד, ומדוד שוב.
\end{enumerate}

\chapter{שזירה}
\label{ch:entangle}

היחידה נסגרת סביב \en{Bell}: הגדרה, מדד, ומה נשבר כשמתעלמים ממנו.

\begin{defbox}[frametitle={הגדרה — שזירה}]
מצב בל הוא שזירה מקסימלית לשני קיוביטים.
\end{defbox}

\begin{equation}
\label{eq:entangle}
\ket{\Phi^+}=\tfrac1{\sqrt2}(\ket{00}+\ket{11})
\end{equation}

המשמעות של~\eqref{eq:entangle}: מודדים את \en{Bell} קודם, ורק אחר כך צוללים לרכיב.

\begin{equation}
\label{eq:entangle-check}
\Delta M \propto \Delta x_{\mathrm{dom}}
\end{equation}
כאן $M$ הוא \en{Bell};~\eqref{eq:entangle-check} מזכירה לחפש רגישות, לא מיקרו־אופטימיזציה.

סדר עבודה:
\begin{enumerate}
\item[(\textbf{א})] מדוד את \en{Bell} לפי~\eqref{eq:entangle}.
\item[(\textbf{ב})] זהה גורם שולט אחד.
\item[(\textbf{ג})] שנה אותו בלבד, ומדוד שוב.
\end{enumerate}

\chapter{אי־שכפול}
\label{ch:noclone}

היחידה נסגרת סביב \en{no-cloning}: הגדרה, מדד, ומה נשבר כשמתעלמים ממנו.

\begin{thmbox}[frametitle={משפט — אי־שכפול}]
אין העתקה אוניטרית של מצב שרירותי.
\end{thmbox}

\begin{equation}
\label{eq:noclone}
\nexists U:\ U(\ket\psi\otimes\ket0)=\ket\psi\otimes\ket\psi
\end{equation}

המשמעות של~\eqref{eq:noclone}: מודדים את \en{no-cloning} קודם, ורק אחר כך צוללים לרכיב.

\begin{tablebox}\centering\small
\begin{tabular}{@{}lll@{}}
\toprule
\textbf{מצב} & \textbf{סדר גודל} & \textbf{פעולה}\\
\midrule
תקין & $\sim 1\times$ & שמור מדידה\\
חשוד & $1.5$–$2\times$ & בודקים תלות\\
שבור & $\ge 2\times$ & תיקון מבני\\
\bottomrule
\end{tabular}
\tabcap{כיול מהיר ל־\en{no-cloning}.}
\end{tablebox}

\begin{tikzpic}
\begin{tikzpicture}[font=\small,node distance=8mm,>={Latex[length=2mm]},
  box/.style={draw,thick,rounded corners,align=center,minimum width=2.4cm,minimum height=1.0cm,fill=cBlueBg}]
  \node[box] (a) {\en{In}};
  \node[box,right=of a] (b) {\en{no-cloni}};
  \node[box,right=of b,fill=cGreenBg] (c) {\en{Out}};
  \draw[->,thick] (a) -- (b);
  \draw[->,thick] (b) -- (c);
\end{tikzpicture}
\end{tikzpic}
\begin{center}\small איור — זרימה ל־\en{no-cloning}\end{center}

סדר עבודה:
\begin{enumerate}
\item[(\textbf{א})] מדוד את \en{no-cloning} לפי~\eqref{eq:noclone}.
\item[(\textbf{ב})] זהה גורם שולט אחד.
\item[(\textbf{ג})] שנה אותו בלבד, ומדוד שוב.
\end{enumerate}

\chapter{טלפורטציה}
\label{ch:teleport}

היחידה נסגרת סביב \en{EPR}: הגדרה, מדד, ומה נשבר כשמתעלמים ממנו.

\begin{keybox}[frametitle={רעיון מפתח — טלפורטציה}]
שתי סיביות קלאסיות + שזירה משחזרות את המצב אצל בוב.
\end{keybox}

\begin{equation}
\label{eq:teleport}
\ket\psi_B = X^{m_2}Z^{m_1}\ket\psi
\end{equation}

המשמעות של~\eqref{eq:teleport}: מודדים את \en{EPR} קודם, ורק אחר כך צוללים לרכיב.

\begin{equation}
\label{eq:teleport-check}
\Delta M \propto \Delta x_{\mathrm{dom}}
\end{equation}
כאן $M$ הוא \en{EPR};~\eqref{eq:teleport-check} מזכירה לחפש רגישות, לא מיקרו־אופטימיזציה.

סדר עבודה:
\begin{enumerate}
\item[(\textbf{א})] מדוד את \en{EPR} לפי~\eqref{eq:teleport}.
\item[(\textbf{ב})] זהה גורם שולט אחד.
\item[(\textbf{ג})] שנה אותו בלבד, ומדוד שוב.
\end{enumerate}

\chapter{אלגוריתם גרובר}
\label{ch:grover}

היחידה נסגרת סביב \en{oracle}: הגדרה, מדד, ומה נשבר כשמתעלמים ממנו.

\begin{defbox}[frametitle={הגדרה — אלגוריתם גרובר}]
חיפוש לא ממוין מאיץ ריבועית מול סריקה קלאסית.
\end{defbox}

\begin{equation}
\label{eq:grover}
O(\sqrt N)\ \mathrm{queries}
\end{equation}

המשמעות של~\eqref{eq:grover}: מודדים את \en{oracle} קודם, ורק אחר כך צוללים לרכיב.

רצף שממחיש לחץ על המדד:
\begin{asmblock}
\begin{lstlisting}[style=asm]
lw   $t0, 0($s0)
add  $t1, $t0, $s1   # depends on t0
\end{lstlisting}
\end{asmblock}
התלות על \code{\$t0} היא עלות ב־\en{oracle}, לא ״באג קוד״.

סדר עבודה:
\begin{enumerate}
\item[(\textbf{א})] מדוד את \en{oracle} לפי~\eqref{eq:grover}.
\item[(\textbf{ב})] זהה גורם שולט אחד.
\item[(\textbf{ג})] שנה אותו בלבד, ומדוד שוב.
\end{enumerate}

\chapter{שור}
\label{ch:shor}

היחידה נסגרת סביב \en{period}: הגדרה, מדד, ומה נשבר כשמתעלמים ממנו.

\begin{keybox}[frametitle={רעיון מפתח — שור}]
מציאת הסדר $r$ מפרקת את $N$ בפולינום.
\end{keybox}

\begin{equation}
\label{eq:shor}
a^r \equiv 1 \pmod N
\end{equation}

המשמעות של~\eqref{eq:shor}: מודדים את \en{period} קודם, ורק אחר כך צוללים לרכיב.

\begin{equation}
\label{eq:shor-check}
\Delta M \propto \Delta x_{\mathrm{dom}}
\end{equation}
כאן $M$ הוא \en{period};~\eqref{eq:shor-check} מזכירה לחפש רגישות, לא מיקרו־אופטימיזציה.

סדר עבודה:
\begin{enumerate}
\item[(\textbf{א})] מדוד את \en{period} לפי~\eqref{eq:shor}.
\item[(\textbf{ב})] זהה גורם שולט אחד.
\item[(\textbf{ג})] שנה אותו בלבד, ומדוד שוב.
\end{enumerate}

\chapter{קודים לתיקון שגיאות}
\label{ch:qec}

היחידה נסגרת סביב \en{surface}: הגדרה, מדד, ומה נשבר כשמתעלמים ממנו.

\begin{defbox}[frametitle={הגדרה — קודים לתיקון שגיאות}]
מרחק הקוד קובע כמה שגיאות אפשר לתקן.
\end{defbox}

\begin{equation}
\label{eq:qec}
d_{\mathrm{code}} \ge 2t+1
\end{equation}

המשמעות של~\eqref{eq:qec}: מודדים את \en{surface} קודם, ורק אחר כך צוללים לרכיב.

\begin{tablebox}\centering\small
\begin{tabular}{@{}lll@{}}
\toprule
\textbf{מצב} & \textbf{סדר גודל} & \textbf{פעולה}\\
\midrule
תקין & $\sim 1\times$ & שמור מדידה\\
חשוד & $1.5$–$2\times$ & בודקים תלות\\
שבור & $\ge 2\times$ & תיקון מבני\\
\bottomrule
\end{tabular}
\tabcap{כיול מהיר ל־\en{surface}.}
\end{tablebox}

סדר עבודה:
\begin{enumerate}
\item[(\textbf{א})] מדוד את \en{surface} לפי~\eqref{eq:qec}.
\item[(\textbf{ב})] זהה גורם שולט אחד.
\item[(\textbf{ג})] שנה אותו בלבד, ומדוד שוב.
\end{enumerate}

\chapter{VQE}
\label{ch:vqe}

היחידה נסגרת סביב \en{ansatz}: הגדרה, מדד, ומה נשבר כשמתעלמים ממנו.

\begin{notebox}[frametitle={הערה — VQE}]
לולאה היברידית: מעגל מפרמטרי + אופטימיזציה קלאסית.
\end{notebox}

\begin{equation}
\label{eq:vqe}
E(\theta)=\bra{\psi(\theta)}H\ket{\psi(\theta)}
\end{equation}

המשמעות של~\eqref{eq:vqe}: מודדים את \en{ansatz} קודם, ורק אחר כך צוללים לרכיב.

\begin{equation}
\label{eq:vqe-check}
\Delta M \propto \Delta x_{\mathrm{dom}}
\end{equation}
כאן $M$ הוא \en{ansatz};~\eqref{eq:vqe-check} מזכירה לחפש רגישות, לא מיקרו־אופטימיזציה.

סדר עבודה:
\begin{enumerate}
\item[(\textbf{א})] מדוד את \en{ansatz} לפי~\eqref{eq:vqe}.
\item[(\textbf{ב})] זהה גורם שולט אחד.
\item[(\textbf{ג})] שנה אותו בלבד, ומדוד שוב.
\end{enumerate}

\chapter{רעש ומדידה}
\label{ch:noise}

היחידה נסגרת סביב \en{depolarizing}: הגדרה, מדד, ומה נשבר כשמתעלמים ממנו.

\begin{warnbox}[frametitle={מלכודת — רעש ומדידה}]
רעש דה־פולריזציה מערבב את המצב עם הזהות — סימולציה אידיאלית מטעה.
\end{warnbox}

\begin{equation}
\label{eq:noise}
\mathcal{E}(\rho)=(1-p)\rho+p\tfrac{I}{2}
\end{equation}

המשמעות של~\eqref{eq:noise}: מודדים את \en{depolarizing} קודם, ורק אחר כך צוללים לרכיב.

\begin{tikzpic}
\begin{tikzpicture}[font=\small,node distance=8mm,>={Latex[length=2mm]},
  box/.style={draw,thick,rounded corners,align=center,minimum width=2.4cm,minimum height=1.0cm,fill=cBlueBg}]
  \node[box] (a) {\en{In}};
  \node[box,right=of a] (b) {\en{depolari}};
  \node[box,right=of b,fill=cGreenBg] (c) {\en{Out}};
  \draw[->,thick] (a) -- (b);
  \draw[->,thick] (b) -- (c);
\end{tikzpicture}
\end{tikzpic}
\begin{center}\small איור — זרימה ל־\en{depolarizing}\end{center}

סדר עבודה:
\begin{enumerate}
\item[(\textbf{א})] מדוד את \en{depolarizing} לפי~\eqref{eq:noise}.
\item[(\textbf{ב})] זהה גורם שולט אחד.
\item[(\textbf{ג})] שנה אותו בלבד, ומדוד שוב.
\end{enumerate}


\chapter{חיבור — חישוב קוונטי}
\label{ch:bridge-qubit}

הפרקים בחלק ״חישוב קוונטי״ חולקים משמעת אחת: מדד, נוסחה, מדידה.
הנוסחאות \\eqref{eq:qubit}, \\eqref{eq:gates}, \\eqref{eq:entangle}, \\eqref{eq:noclone} אינן אוסף עובדות — הן חוזה מדידה.

\begin{keybox}[frametitle={חוזה החלק}]
לכל נוסחה ממוספרת בחלק כתוב ניסוי אחד שמאמת אותה. בלי ניסוי — היא דקורציה.
\end{keybox}

כששני מדדים מתנגשים, בוחרים לפי מטרת המערכת (השהיה מול מעבר־נתונים, דיוק מול עלות),
לא לפי הנוסחה היפה יותר.

\part{אבטחה וקריפטוגרפיה}

\chapter{הצפנה סימטרית}
\label{ch:sym}

היחידה נסגרת סביב \en{AES}: הגדרה, מדד, ומה נשבר כשמתעלמים ממנו.

\begin{defbox}[frametitle={הגדרה — הצפנה סימטרית}]
מפתח אחד מצפין ומפענח; סודיות המפתח היא הכל.
\end{defbox}

\begin{equation}
\label{eq:sym}
C = E_K(P)
\end{equation}

המשמעות של~\eqref{eq:sym}: מודדים את \en{AES} קודם, ורק אחר כך צוללים לרכיב.

\begin{equation}
\label{eq:sym-check}
\Delta M \propto \Delta x_{\mathrm{dom}}
\end{equation}
כאן $M$ הוא \en{AES};~\eqref{eq:sym-check} מזכירה לחפש רגישות, לא מיקרו־אופטימיזציה.

רצף שממחיש לחץ על המדד:
\begin{asmblock}
\begin{lstlisting}[style=asm]
lw   $t0, 0($s0)
add  $t1, $t0, $s1   # depends on t0
\end{lstlisting}
\end{asmblock}
התלות על \code{\$t0} היא עלות ב־\en{AES}, לא ״באג קוד״.

סדר עבודה:
\begin{enumerate}
\item[(\textbf{א})] מדוד את \en{AES} לפי~\eqref{eq:sym}.
\item[(\textbf{ב})] זהה גורם שולט אחד.
\item[(\textbf{ג})] שנה אותו בלבד, ומדוד שוב.
\end{enumerate}

\chapter{מפתח ציבורי}
\label{ch:pubkey}

היחידה נסגרת סביב \en{RSA}: הגדרה, מדד, ומה נשבר כשמתעלמים ממנו.

\begin{keybox}[frametitle={רעיון מפתח — מפתח ציבורי}]
הצפנה במפתח ציבורי; פיענוח רק למי שמחזיק $d$.
\end{keybox}

\begin{equation}
\label{eq:pubkey}
c \equiv m^e \pmod n
\end{equation}

המשמעות של~\eqref{eq:pubkey}: מודדים את \en{RSA} קודם, ורק אחר כך צוללים לרכיב.

\begin{tablebox}\centering\small
\begin{tabular}{@{}lll@{}}
\toprule
\textbf{מצב} & \textbf{סדר גודל} & \textbf{פעולה}\\
\midrule
תקין & $\sim 1\times$ & שמור מדידה\\
חשוד & $1.5$–$2\times$ & בודקים תלות\\
שבור & $\ge 2\times$ & תיקון מבני\\
\bottomrule
\end{tabular}
\tabcap{כיול מהיר ל־\en{RSA}.}
\end{tablebox}

סדר עבודה:
\begin{enumerate}
\item[(\textbf{א})] מדוד את \en{RSA} לפי~\eqref{eq:pubkey}.
\item[(\textbf{ב})] זהה גורם שולט אחד.
\item[(\textbf{ג})] שנה אותו בלבד, ומדוד שוב.
\end{enumerate}

\chapter{חתימות}
\label{ch:sig}

היחידה נסגרת סביב \en{verify}: הגדרה, מדד, ומה נשבר כשמתעלמים ממנו.

\begin{defbox}[frametitle={הגדרה — חתימות}]
חתימה מאמתת מקור ושלמות בלי לחשוף מפתח פרטי.
\end{defbox}

\begin{equation}
\label{eq:sig}
\mathsf{Vrfy}(pk,\sigma,m)=1
\end{equation}

המשמעות של~\eqref{eq:sig}: מודדים את \en{verify} קודם, ורק אחר כך צוללים לרכיב.

\begin{equation}
\label{eq:sig-check}
\Delta M \propto \Delta x_{\mathrm{dom}}
\end{equation}
כאן $M$ הוא \en{verify};~\eqref{eq:sig-check} מזכירה לחפש רגישות, לא מיקרו־אופטימיזציה.

סדר עבודה:
\begin{enumerate}
\item[(\textbf{א})] מדוד את \en{verify} לפי~\eqref{eq:sig}.
\item[(\textbf{ב})] זהה גורם שולט אחד.
\item[(\textbf{ג})] שנה אותו בלבד, ומדוד שוב.
\end{enumerate}

\chapter{גיבוב קריפטוגרפי}
\label{ch:hashc}

היחידה נסגרת סביב \en{SHA}: הגדרה, מדד, ומה נשבר כשמתעלמים ממנו.

\begin{keybox}[frametitle={רעיון מפתח — גיבוב קריפטוגרפי}]
עמידות להתנגשות חשובה יותר מאורך הפלט לבדו.
\end{keybox}

\begin{equation}
\label{eq:hashc}
H:\{0,1\}^*\to\{0,1\}^n
\end{equation}

המשמעות של~\eqref{eq:hashc}: מודדים את \en{SHA} קודם, ורק אחר כך צוללים לרכיב.

סדר עבודה:
\begin{enumerate}
\item[(\textbf{א})] מדוד את \en{SHA} לפי~\eqref{eq:hashc}.
\item[(\textbf{ב})] זהה גורם שולט אחד.
\item[(\textbf{ג})] שנה אותו בלבד, ומדוד שוב.
\end{enumerate}

\chapter{פרוטוקולי אימות}
\label{ch:auth}

היחידה נסגרת סביב \en{challenge}: הגדרה, מדד, ומה נשבר כשמתעלמים ממנו.

\begin{thmbox}[frametitle={משפט — פרוטוקולי אימות}]
פרוטוקול בטוח אם יתרון היריב זניח.
\end{thmbox}

\begin{equation}
\label{eq:auth}
\mathsf{Adv} \le \varepsilon
\end{equation}

המשמעות של~\eqref{eq:auth}: מודדים את \en{challenge} קודם, ורק אחר כך צוללים לרכיב.

\begin{equation}
\label{eq:auth-check}
\Delta M \propto \Delta x_{\mathrm{dom}}
\end{equation}
כאן $M$ הוא \en{challenge};~\eqref{eq:auth-check} מזכירה לחפש רגישות, לא מיקרו־אופטימיזציה.

סדר עבודה:
\begin{enumerate}
\item[(\textbf{א})] מדוד את \en{challenge} לפי~\eqref{eq:auth}.
\item[(\textbf{ב})] זהה גורם שולט אחד.
\item[(\textbf{ג})] שנה אותו בלבד, ומדוד שוב.
\end{enumerate}

\chapter{ערוץ מאובטח}
\label{ch:securech}

היחידה נסגרת סביב \en{AEAD}: הגדרה, מדד, ומה נשבר כשמתעלמים ממנו.

\begin{defbox}[frametitle={הגדרה — ערוץ מאובטח}]
\en{AEAD} מצפין ומאמת יחד — כולל נתונים נלווים.
\end{defbox}

\begin{equation}
\label{eq:securech}
\mathsf{Enc}_k(\mathrm{nonce},\mathrm{aad},\mathrm{pt})
\end{equation}

המשמעות של~\eqref{eq:securech}: מודדים את \en{AEAD} קודם, ורק אחר כך צוללים לרכיב.

\begin{tablebox}\centering\small
\begin{tabular}{@{}lll@{}}
\toprule
\textbf{מצב} & \textbf{סדר גודל} & \textbf{פעולה}\\
\midrule
תקין & $\sim 1\times$ & שמור מדידה\\
חשוד & $1.5$–$2\times$ & בודקים תלות\\
שבור & $\ge 2\times$ & תיקון מבני\\
\bottomrule
\end{tabular}
\tabcap{כיול מהיר ל־\en{AEAD}.}
\end{tablebox}

רצף שממחיש לחץ על המדד:
\begin{asmblock}
\begin{lstlisting}[style=asm]
lw   $t0, 0($s0)
add  $t1, $t0, $s1   # depends on t0
\end{lstlisting}
\end{asmblock}
התלות על \code{\$t0} היא עלות ב־\en{AEAD}, לא ״באג קוד״.

\begin{tikzpic}
\begin{tikzpicture}[font=\small,node distance=8mm,>={Latex[length=2mm]},
  box/.style={draw,thick,rounded corners,align=center,minimum width=2.4cm,minimum height=1.0cm,fill=cBlueBg}]
  \node[box] (a) {\en{In}};
  \node[box,right=of a] (b) {\en{AEAD}};
  \node[box,right=of b,fill=cGreenBg] (c) {\en{Out}};
  \draw[->,thick] (a) -- (b);
  \draw[->,thick] (b) -- (c);
\end{tikzpicture}
\end{tikzpic}
\begin{center}\small איור — זרימה ל־\en{AEAD}\end{center}

סדר עבודה:
\begin{enumerate}
\item[(\textbf{א})] מדוד את \en{AEAD} לפי~\eqref{eq:securech}.
\item[(\textbf{ב})] זהה גורם שולט אחד.
\item[(\textbf{ג})] שנה אותו בלבד, ומדוד שוב.
\end{enumerate}


\chapter{חיבור — אבטחה וקריפטוגרפיה}
\label{ch:bridge-sym}

הפרקים בחלק ״אבטחה וקריפטוגרפיה״ חולקים משמעת אחת: מדד, נוסחה, מדידה.
הנוסחאות \\eqref{eq:sym}, \\eqref{eq:pubkey}, \\eqref{eq:sig}, \\eqref{eq:hashc} אינן אוסף עובדות — הן חוזה מדידה.

\begin{keybox}[frametitle={חוזה החלק}]
לכל נוסחה ממוספרת בחלק כתוב ניסוי אחד שמאמת אותה. בלי ניסוי — היא דקורציה.
\end{keybox}

כששני מדדים מתנגשים, בוחרים לפי מטרת המערכת (השהיה מול מעבר־נתונים, דיוק מול עלות),
לא לפי הנוסחה היפה יותר.


\part{תבניות עבודה}
\chapter{מעגל בל ואי־שכפול}
\label{ch:bell-template}

מצב בל הוא תבנית העבודה לשזירה דו־קיוביטית.

\begin{equation}
\label{eq:bell-template}
\ket{\Phi^+} = \tfrac1{\sqrt2}(\ket{00}+\ket{11})
\end{equation}

\begin{qcirc}
\begin{quantikz}
\lstick{$\ket0$} & \gate{H} & \ctrl{1} & \meter{} \\
\lstick{$\ket0$} & \qw & \targ{} & \meter{}
\end{quantikz}
\end{qcirc}
\begin{center}\small איור — הכנת $\ket{\Phi^+}$ ומדידה\end{center}

\begin{thmbox}[frametitle={משפט — אי־שכפול}]
לא קיים אופרטור אוניטרי $U$ כך ש־$U(\ket\psi\otimes\ket0)=\ket\psi\otimes\ket\psi$ לכל $\ket\psi$.
\end{thmbox}

\noindent\textbf{הוכחה (בקצרה).}
 נניח שקיים $U$ לשני מצבים. אז $\braket{\psi}{\phi}=\braket{\psi}{\phi}^2$,
ולכן $\braket{\psi}{\phi}\in\{0,1\}$ — סתירה לליניאריות על מצב שרירותי. $\qquad\blacksquare$

הלקח התפעולי: שזירה לא ניתנת להעתקה מקומית; כל פרוטוקול שמניח עותק שקט של מצב שזור שבור מראש.


\chapter{עקומות ביצוע}
\label{ch:curves}

גרף הוא מדידה: ציר $x$ פרמטר תכנון, ציר $y$ המדד.

\begin{tikzpic}
\begin{tikzpicture}
\begin{axis}[
  width=11cm, height=6cm,
  xlabel={גודל (יח')},
  ylabel={עלות / שגיאה},
  legend style={at={(0.97,0.97)},anchor=north east,font=\small},
  grid=major, grid style={gray!20},
]
\addplot[thick,mark=*] coordinates {(1,18)(2,14)(4,10)(8,7)(16,5)(32,4)};
\addlegendentry{בסיס}
\addplot[thick,dashed,mark=square*] coordinates {(1,14)(2,10)(4,7)(8,4.5)(16,3)(32,2.2)};
\addlegendentry{משופר}
\end{axis}
\end{tikzpicture}
\end{tikzpic}
\begin{center}\small איור — ירידת שגיאה עם גודל המשאב\end{center}

נקודת התכנון היא שם שהתועלת השולית יורדת מתחת לעלות השולית — לא קצה העקומה.

\part{תרגילים פתורים}
\chapter{תרגילים מסכמים}
\label{ch:ex}


\begin{exbox}[frametitle={תרגיל 1 — \en{CPI}}]
\textbf{(א)} חשב לפי~\eqref{eq:pipeline} עם $a=1$, $b=0.05$, $c=80$.\\[3pt]
\textbf{(ב)} הורד את $b$ ל־$0.02$ והעלה את $a$ ל־$2$. האם כדאי?\\[3pt]
\textbf{(ג)} מתי הקירוב נשבר?
\end{exbox}

\begin{keybox}[frametitle={פתרון 1},nobreak=true]
\textbf{פתרון (א).} מצבים ל~\eqref{eq:pipeline} ומקבלים ערך יחיד.\\[2pt]
\textbf{פתרון (ב).} משווים שני ערכים של המדד — ההחלטה לפי המספר, לא לפי אינטואיציה.\\[2pt]
\textbf{פתרון (ג).} כשגורמים שנחשבו בלתי־תלויים מתואמים.\\[2pt]
\textbf{לקח.} מדוד דרך הנוסחה.
\end{keybox}


\begin{exbox}[frametitle={תרגיל 2 — \en{bypass}}]
\textbf{(א)} חשב לפי~\eqref{eq:forward} עם $a=1$, $b=0.05$, $c=80$.\\[3pt]
\textbf{(ב)} הורד את $b$ ל־$0.02$ והעלה את $a$ ל־$2$. האם כדאי?\\[3pt]
\textbf{(ג)} מתי הקירוב נשבר?
\end{exbox}

\begin{keybox}[frametitle={פתרון 2},nobreak=true]
\textbf{פתרון (א).} מצבים ל~\eqref{eq:forward} ומקבלים ערך יחיד.\\[2pt]
\textbf{פתרון (ב).} משווים שני ערכים של המדד — ההחלטה לפי המספר, לא לפי אינטואיציה.\\[2pt]
\textbf{פתרון (ג).} כשגורמים שנחשבו בלתי־תלויים מתואמים.\\[2pt]
\textbf{לקח.} מדוד דרך הנוסחה.
\end{keybox}


\begin{exbox}[frametitle={תרגיל 3 — \en{ROB}}]
\textbf{(א)} חשב לפי~\eqref{eq:ooo} עם $a=1$, $b=0.05$, $c=80$.\\[3pt]
\textbf{(ב)} הורד את $b$ ל־$0.02$ והעלה את $a$ ל־$2$. האם כדאי?\\[3pt]
\textbf{(ג)} מתי הקירוב נשבר?
\end{exbox}

\begin{keybox}[frametitle={פתרון 3},nobreak=true]
\textbf{פתרון (א).} מצבים ל~\eqref{eq:ooo} ומקבלים ערך יחיד.\\[2pt]
\textbf{פתרון (ב).} משווים שני ערכים של המדד — ההחלטה לפי המספר, לא לפי אינטואיציה.\\[2pt]
\textbf{פתרון (ג).} כשגורמים שנחשבו בלתי־תלויים מתואמים.\\[2pt]
\textbf{לקח.} מדוד דרך הנוסחה.
\end{keybox}


\begin{exbox}[frametitle={תרגיל 4 — \en{inclusive}}]
\textbf{(א)} חשב לפי~\eqref{eq:l2} עם $a=1$, $b=0.05$, $c=80$.\\[3pt]
\textbf{(ב)} הורד את $b$ ל־$0.02$ והעלה את $a$ ל־$2$. האם כדאי?\\[3pt]
\textbf{(ג)} מתי הקירוב נשבר?
\end{exbox}

\begin{keybox}[frametitle={פתרון 4},nobreak=true]
\textbf{פתרון (א).} מצבים ל~\eqref{eq:l2} ומקבלים ערך יחיד.\\[2pt]
\textbf{פתרון (ב).} משווים שני ערכים של המדד — ההחלטה לפי המספר, לא לפי אינטואיציה.\\[2pt]
\textbf{פתרון (ג).} כשגורמים שנחשבו בלתי־תלויים מתואמים.\\[2pt]
\textbf{לקח.} מדוד דרך הנוסחה.
\end{keybox}


\begin{exbox}[frametitle={תרגיל 5 — \en{VPN}}]
\textbf{(א)} חשב לפי~\eqref{eq:tlb} עם $a=1$, $b=0.05$, $c=80$.\\[3pt]
\textbf{(ב)} הורד את $b$ ל־$0.02$ והעלה את $a$ ל־$2$. האם כדאי?\\[3pt]
\textbf{(ג)} מתי הקירוב נשבר?
\end{exbox}

\begin{keybox}[frametitle={פתרון 5},nobreak=true]
\textbf{פתרון (א).} מצבים ל~\eqref{eq:tlb} ומקבלים ערך יחיד.\\[2pt]
\textbf{פתרון (ב).} משווים שני ערכים של המדד — ההחלטה לפי המספר, לא לפי אינטואיציה.\\[2pt]
\textbf{פתרון (ג).} כשגורמים שנחשבו בלתי־תלויים מתואמים.\\[2pt]
\textbf{לקח.} מדוד דרך הנוסחה.
\end{keybox}


\begin{exbox}[frametitle={תרגיל 6 — \en{bandwidth}}]
\textbf{(א)} חשב לפי~\eqref{eq:membw} עם $a=1$, $b=0.05$, $c=80$.\\[3pt]
\textbf{(ב)} הורד את $b$ ל־$0.02$ והעלה את $a$ ל־$2$. האם כדאי?\\[3pt]
\textbf{(ג)} מתי הקירוב נשבר?
\end{exbox}

\begin{keybox}[frametitle={פתרון 6},nobreak=true]
\textbf{פתרון (א).} מצבים ל~\eqref{eq:membw} ומקבלים ערך יחיד.\\[2pt]
\textbf{פתרון (ב).} משווים שני ערכים של המדד — ההחלטה לפי המספר, לא לפי אינטואיציה.\\[2pt]
\textbf{פתרון (ג).} כשגורמים שנחשבו בלתי־תלויים מתואמים.\\[2pt]
\textbf{לקח.} מדוד דרך הנוסחה.
\end{keybox}


\begin{exbox}[frametitle={תרגיל 7 — \en{mutex}}]
\textbf{(א)} חשב לפי~\eqref{eq:sync} עם $a=1$, $b=0.05$, $c=80$.\\[3pt]
\textbf{(ב)} הורד את $b$ ל־$0.02$ והעלה את $a$ ל־$2$. האם כדאי?\\[3pt]
\textbf{(ג)} מתי הקירוב נשבר?
\end{exbox}

\begin{keybox}[frametitle={פתרון 7},nobreak=true]
\textbf{פתרון (א).} מצבים ל~\eqref{eq:sync} ומקבלים ערך יחיד.\\[2pt]
\textbf{פתרון (ב).} משווים שני ערכים של המדד — ההחלטה לפי המספר, לא לפי אינטואיציה.\\[2pt]
\textbf{פתרון (ג).} כשגורמים שנחשבו בלתי־תלויים מתואמים.\\[2pt]
\textbf{לקח.} מדוד דרך הנוסחה.
\end{keybox}


\begin{exbox}[frametitle={תרגיל 8 — \en{DMA}}]
\textbf{(א)} חשב לפי~\eqref{eq:io} עם $a=1$, $b=0.05$, $c=80$.\\[3pt]
\textbf{(ב)} הורד את $b$ ל־$0.02$ והעלה את $a$ ל־$2$. האם כדאי?\\[3pt]
\textbf{(ג)} מתי הקירוב נשבר?
\end{exbox}

\begin{keybox}[frametitle={פתרון 8},nobreak=true]
\textbf{פתרון (א).} מצבים ל~\eqref{eq:io} ומקבלים ערך יחיד.\\[2pt]
\textbf{פתרון (ב).} משווים שני ערכים של המדד — ההחלטה לפי המספר, לא לפי אינטואיציה.\\[2pt]
\textbf{פתרון (ג).} כשגורמים שנחשבו בלתי־תלויים מתואמים.\\[2pt]
\textbf{לקח.} מדוד דרך הנוסחה.
\end{keybox}


\begin{exbox}[frametitle={תרגיל 9 — \en{OSI}}]
\textbf{(א)} חשב לפי~\eqref{eq:layers} עם $a=1$, $b=0.05$, $c=80$.\\[3pt]
\textbf{(ב)} הורד את $b$ ל־$0.02$ והעלה את $a$ ל־$2$. האם כדאי?\\[3pt]
\textbf{(ג)} מתי הקירוב נשבר?
\end{exbox}

\begin{keybox}[frametitle={פתרון 9},nobreak=true]
\textbf{פתרון (א).} מצבים ל~\eqref{eq:layers} ומקבלים ערך יחיד.\\[2pt]
\textbf{פתרון (ב).} משווים שני ערכים של המדד — ההחלטה לפי המספר, לא לפי אינטואיציה.\\[2pt]
\textbf{פתרון (ג).} כשגורמים שנחשבו בלתי־תלויים מתואמים.\\[2pt]
\textbf{לקח.} מדוד דרך הנוסחה.
\end{keybox}


\begin{exbox}[frametitle={תרגיל 10 — \en{latency}}]
\textbf{(א)} חשב לפי~\eqref{eq:qos} עם $a=1$, $b=0.05$, $c=80$.\\[3pt]
\textbf{(ב)} הורד את $b$ ל־$0.02$ והעלה את $a$ ל־$2$. האם כדאי?\\[3pt]
\textbf{(ג)} מתי הקירוב נשבר?
\end{exbox}

\begin{keybox}[frametitle={פתרון 10},nobreak=true]
\textbf{פתרון (א).} מצבים ל~\eqref{eq:qos} ומקבלים ערך יחיד.\\[2pt]
\textbf{פתרון (ב).} משווים שני ערכים של המדד — ההחלטה לפי המספר, לא לפי אינטואיציה.\\[2pt]
\textbf{פתרון (ג).} כשגורמים שנחשבו בלתי־תלויים מתואמים.\\[2pt]
\textbf{לקח.} מדוד דרך הנוסחה.
\end{keybox}


\begin{exbox}[frametitle={תרגיל 11 — \en{BigO}}]
\textbf{(א)} חשב לפי~\eqref{eq:complexity} עם $a=1$, $b=0.05$, $c=80$.\\[3pt]
\textbf{(ב)} הורד את $b$ ל־$0.02$ והעלה את $a$ ל־$2$. האם כדאי?\\[3pt]
\textbf{(ג)} מתי הקירוב נשבר?
\end{exbox}

\begin{keybox}[frametitle={פתרון 11},nobreak=true]
\textbf{פתרון (א).} מצבים ל~\eqref{eq:complexity} ומקבלים ערך יחיד.\\[2pt]
\textbf{פתרון (ב).} משווים שני ערכים של המדד — ההחלטה לפי המספר, לא לפי אינטואיציה.\\[2pt]
\textbf{פתרון (ג).} כשגורמים שנחשבו בלתי־תלויים מתואמים.\\[2pt]
\textbf{לקח.} מדוד דרך הנוסחה.
\end{keybox}


\begin{exbox}[frametitle={תרגיל 12 — \en{opt}}]
\textbf{(א)} חשב לפי~\eqref{eq:dp} עם $a=1$, $b=0.05$, $c=80$.\\[3pt]
\textbf{(ב)} הורד את $b$ ל־$0.02$ והעלה את $a$ ל־$2$. האם כדאי?\\[3pt]
\textbf{(ג)} מתי הקירוב נשבר?
\end{exbox}

\begin{keybox}[frametitle={פתרון 12},nobreak=true]
\textbf{פתרון (א).} מצבים ל~\eqref{eq:dp} ומקבלים ערך יחיד.\\[2pt]
\textbf{פתרון (ב).} משווים שני ערכים של המדד — ההחלטה לפי המספר, לא לפי אינטואיציה.\\[2pt]
\textbf{פתרון (ג).} כשגורמים שנחשבו בלתי־תלויים מתואמים.\\[2pt]
\textbf{לקח.} מדוד דרך הנוסחה.
\end{keybox}


\begin{exbox}[frametitle={תרגיל 13 — \en{maxflow}}]
\textbf{(א)} חשב לפי~\eqref{eq:flow} עם $a=1$, $b=0.05$, $c=80$.\\[3pt]
\textbf{(ב)} הורד את $b$ ל־$0.02$ והעלה את $a$ ל־$2$. האם כדאי?\\[3pt]
\textbf{(ג)} מתי הקירוב נשבר?
\end{exbox}

\begin{keybox}[frametitle={פתרון 13},nobreak=true]
\textbf{פתרון (א).} מצבים ל~\eqref{eq:flow} ומקבלים ערך יחיד.\\[2pt]
\textbf{פתרון (ב).} משווים שני ערכים של המדד — ההחלטה לפי המספר, לא לפי אינטואיציה.\\[2pt]
\textbf{פתרון (ג).} כשגורמים שנחשבו בלתי־תלויים מתואמים.\\[2pt]
\textbf{לקח.} מדוד דרך הנוסחה.
\end{keybox}


\begin{exbox}[frametitle={תרגיל 14 — \en{logistic}}]
\textbf{(א)} חשב לפי~\eqref{eq:clf} עם $a=1$, $b=0.05$, $c=80$.\\[3pt]
\textbf{(ב)} הורד את $b$ ל־$0.02$ והעלה את $a$ ל־$2$. האם כדאי?\\[3pt]
\textbf{(ג)} מתי הקירוב נשבר?
\end{exbox}

\begin{keybox}[frametitle={פתרון 14},nobreak=true]
\textbf{פתרון (א).} מצבים ל~\eqref{eq:clf} ומקבלים ערך יחיד.\\[2pt]
\textbf{פתרון (ב).} משווים שני ערכים של המדד — ההחלטה לפי המספר, לא לפי אינטואיציה.\\[2pt]
\textbf{פתרון (ג).} כשגורמים שנחשבו בלתי־תלויים מתואמים.\\[2pt]
\textbf{לקח.} מדוד דרך הנוסחה.
\end{keybox}


\begin{exbox}[frametitle={תרגיל 15 — \en{VC}}]
\textbf{(א)} חשב לפי~\eqref{eq:gen} עם $a=1$, $b=0.05$, $c=80$.\\[3pt]
\textbf{(ב)} הורד את $b$ ל־$0.02$ והעלה את $a$ ל־$2$. האם כדאי?\\[3pt]
\textbf{(ג)} מתי הקירוב נשבר?
\end{exbox}

\begin{keybox}[frametitle={פתרון 15},nobreak=true]
\textbf{פתרון (א).} מצבים ל~\eqref{eq:gen} ומקבלים ערך יחיד.\\[2pt]
\textbf{פתרון (ב).} משווים שני ערכים של המדד — ההחלטה לפי המספר, לא לפי אינטואיציה.\\[2pt]
\textbf{פתרון (ג).} כשגורמים שנחשבו בלתי־תלויים מתואמים.\\[2pt]
\textbf{לקח.} מדוד דרך הנוסחה.
\end{keybox}


\begin{exbox}[frametitle={תרגיל 16 — \en{Bloch}}]
\textbf{(א)} חשב לפי~\eqref{eq:qubit} עם $a=1$, $b=0.05$, $c=80$.\\[3pt]
\textbf{(ב)} הורד את $b$ ל־$0.02$ והעלה את $a$ ל־$2$. האם כדאי?\\[3pt]
\textbf{(ג)} מתי הקירוב נשבר?
\end{exbox}

\begin{keybox}[frametitle={פתרון 16},nobreak=true]
\textbf{פתרון (א).} מצבים ל~\eqref{eq:qubit} ומקבלים ערך יחיד.\\[2pt]
\textbf{פתרון (ב).} משווים שני ערכים של המדד — ההחלטה לפי המספר, לא לפי אינטואיציה.\\[2pt]
\textbf{פתרון (ג).} כשגורמים שנחשבו בלתי־תלויים מתואמים.\\[2pt]
\textbf{לקח.} מדוד דרך הנוסחה.
\end{keybox}


\begin{exbox}[frametitle={תרגיל 17 — \en{no-cloning}}]
\textbf{(א)} חשב לפי~\eqref{eq:noclone} עם $a=1$, $b=0.05$, $c=80$.\\[3pt]
\textbf{(ב)} הורד את $b$ ל־$0.02$ והעלה את $a$ ל־$2$. האם כדאי?\\[3pt]
\textbf{(ג)} מתי הקירוב נשבר?
\end{exbox}

\begin{keybox}[frametitle={פתרון 17},nobreak=true]
\textbf{פתרון (א).} מצבים ל~\eqref{eq:noclone} ומקבלים ערך יחיד.\\[2pt]
\textbf{פתרון (ב).} משווים שני ערכים של המדד — ההחלטה לפי המספר, לא לפי אינטואיציה.\\[2pt]
\textbf{פתרון (ג).} כשגורמים שנחשבו בלתי־תלויים מתואמים.\\[2pt]
\textbf{לקח.} מדוד דרך הנוסחה.
\end{keybox}


\begin{exbox}[frametitle={תרגיל 18 — \en{period}}]
\textbf{(א)} חשב לפי~\eqref{eq:shor} עם $a=1$, $b=0.05$, $c=80$.\\[3pt]
\textbf{(ב)} הורד את $b$ ל־$0.02$ והעלה את $a$ ל־$2$. האם כדאי?\\[3pt]
\textbf{(ג)} מתי הקירוב נשבר?
\end{exbox}

\begin{keybox}[frametitle={פתרון 18},nobreak=true]
\textbf{פתרון (א).} מצבים ל~\eqref{eq:shor} ומקבלים ערך יחיד.\\[2pt]
\textbf{פתרון (ב).} משווים שני ערכים של המדד — ההחלטה לפי המספר, לא לפי אינטואיציה.\\[2pt]
\textbf{פתרון (ג).} כשגורמים שנחשבו בלתי־תלויים מתואמים.\\[2pt]
\textbf{לקח.} מדוד דרך הנוסחה.
\end{keybox}


\begin{exbox}[frametitle={תרגיל 19 — \en{depolarizing}}]
\textbf{(א)} חשב לפי~\eqref{eq:noise} עם $a=1$, $b=0.05$, $c=80$.\\[3pt]
\textbf{(ב)} הורד את $b$ ל־$0.02$ והעלה את $a$ ל־$2$. האם כדאי?\\[3pt]
\textbf{(ג)} מתי הקירוב נשבר?
\end{exbox}

\begin{keybox}[frametitle={פתרון 19},nobreak=true]
\textbf{פתרון (א).} מצבים ל~\eqref{eq:noise} ומקבלים ערך יחיד.\\[2pt]
\textbf{פתרון (ב).} משווים שני ערכים של המדד — ההחלטה לפי המספר, לא לפי אינטואיציה.\\[2pt]
\textbf{פתרון (ג).} כשגורמים שנחשבו בלתי־תלויים מתואמים.\\[2pt]
\textbf{לקח.} מדוד דרך הנוסחה.
\end{keybox}


\begin{exbox}[frametitle={תרגיל 20 — \en{verify}}]
\textbf{(א)} חשב לפי~\eqref{eq:sig} עם $a=1$, $b=0.05$, $c=80$.\\[3pt]
\textbf{(ב)} הורד את $b$ ל־$0.02$ והעלה את $a$ ל־$2$. האם כדאי?\\[3pt]
\textbf{(ג)} מתי הקירוב נשבר?
\end{exbox}

\begin{keybox}[frametitle={פתרון 20},nobreak=true]
\textbf{פתרון (א).} מצבים ל~\eqref{eq:sig} ומקבלים ערך יחיד.\\[2pt]
\textbf{פתרון (ב).} משווים שני ערכים של המדד — ההחלטה לפי המספר, לא לפי אינטואיציה.\\[2pt]
\textbf{פתרון (ג).} כשגורמים שנחשבו בלתי־תלויים מתואמים.\\[2pt]
\textbf{לקח.} מדוד דרך הנוסחה.
\end{keybox}


\chapter{נספח — זהויות}
\label{ch:appendix}

\[
T = \mathrm{IC}\times\mathrm{CPI}\times T_{\mathrm{clk}},\qquad
T_{\mathrm{avg}}=T_{\mathrm{hit}}+r_{\mathrm{miss}}T_{\mathrm{miss}}
\]
\[
HXH=Z,\qquad HZH=X,\qquad H^2=I
\]

\begin{tablebox}\centering\small
\begin{tabular}{@{}ll@{}}
\toprule
\textbf{קיצור} & \textbf{משמעות}\\
\midrule
\en{CPI} & מחזורים להוראה\\
\en{AMAT} & זמן גישה ממוצע\\
\en{IPC} & הוראות למחזור\\
\en{TLB} & מטמון תרגום\\
\en{VQE} & אומדן ערך עצמי וריאציוני\\
\bottomrule
\end{tabular}
\tabcap{קיצורים מרכזיים.}
\end{tablebox}

\end{document}
